\documentclass[11pt]{article}
\usepackage[letterpaper,margin=1in]{geometry}
\usepackage{amsmath,amssymb,amsthm,mathtools}
\usepackage{booktabs,tabularx,array,microtype,xurl,longtable}
\usepackage[colorlinks=true,linkcolor=blue,citecolor=blue,urlcolor=blue]{hyperref}
\renewcommand{\labelitemi}{\ensuremath{\bullet}}
\newtheorem{theorem}{Theorem}[section]
\newtheorem{proposition}[theorem]{Proposition}
\newtheorem{lemma}[theorem]{Lemma}
\newtheorem{corollary}[theorem]{Corollary}
\theoremstyle{definition}
\newtheorem{definition}[theorem]{Definition}
\theoremstyle{remark}
\newtheorem{remark}[theorem]{Remark}
\newcommand{\R}{\mathbb R}
\newcommand{\C}{\mathbb C}
\newcommand{\E}{\mathbb E}
\newcommand{\dd}{\,\mathrm d}
\newcommand{\cG}{\mathcal G}
\newcommand{\cL}{\mathcal L}
\newcommand{\fa}{\mathfrak a}
\newcommand{\Rem}[1]{[\![#1]\!]}
\newcommand{\re}{\operatorname{Re}}
\newcommand{\im}{\operatorname{Im}}
\numberwithin{equation}{section}
\title{Height-$p$ transport for heat-flowed $L$-functions\\[2pt]
\large An explicit remainder for the exterior input of strip comparisons}
\author{Drafted for Edward Baker}
\date{17 October 2026\\\small Unreviewed working draft}
\begin{document}
\maketitle
\begin{abstract}
For the de Bruijn--Newman heat flow $H_t$ of the Riemann $\xi$-function and
of the completed Dirichlet $L$-functions $\xi_\chi$, strip comparisons
bound the squared height of the highest zero by a differential inequality
whose coefficient is $S(x,p,t)=-\im(H_t'/H_t)(x+ip)$, read at a fixed
height $p$ above the zeros. We show that $S$ is half the real part of the
logarithmic derivative of a Gaussian average of the completed function,
and that recentring this average at the saddle point of the Gamma factor
gives an exact transport identity. We bound its remainder explicitly
using only Euler-product majorants at real arguments, the digamma function
and a partial-summation bound in the critical strip, with no
Riemann--Siegel formula, truncated Dirichlet series or zero-free region.
Cauchy's estimate then bounds $\partial_pS$. Evaluated in floating point
on a grid of times, the bound keeps the exact time-zero floor of $S$ with
no loss for every primitive character of conductor at least
$1.1\cdot10^5$ (even) or $4.9\cdot10^4$ (odd) at $p=\sqrt5$, and for $\zeta$
beyond the heights of the published computations of Polymath and of Platt
and Trudgian. At Hakimoglu-Brown's height $p=61/20$ its remainder,
evaluated in floating point, is $1.7$ to $4.4$ times smaller than that of
his Riemann--Siegel-based estimate. As
a cross-check we restart his derivative-kernel comparison from the
classical strip. A 34-row table, verified in exact and directed-rounding
arithmetic, gives $\Lambda\le0.18891$ conditionally on the analytic lemmas
of his unrefereed preprint. This is not a record: unrefereed preprints
already claim stronger bounds from the same starting point.
\end{abstract}

\noindent\textbf{Draft status.}
This working manuscript was requested before the formal review of the
project's first round. It integrates round-one Notes~01 and~02
\cite{n01,n02} and their check programmes \cite{records}. The analytic
statements of Sections~\ref{sec:setting}--\ref{sec:derivative} are proved
here. Section~\ref{sec:numerics} reports floating-point evaluations of
those bounds, not outward-rounded certificates. Proposition~\ref{prop:table}
is conditional on a cited unrefereed preprint, and
Section~\ref{sec:planning} is planning arithmetic. Derivations, programmes
and authoring checks used Claude Opus 5.5 (Anthropic). Those checks were
internal and same-model; cross-model and specialist human review are
pending.

\section{Introduction}\label{sec:intro}

\subsection{The heat flow and the strip comparison}

Let $\Phi$ be the Riemann kernel and
$H_t(z)=\int_0^\infty e^{tu^2}\Phi(u)\cos(zu)\dd u$, so that
$H_0(z)=\frac18\xi\bigl(\frac{1+iz}2\bigr)$ and $\partial_tH_t=-\partial_z^2H_t$
\cite{polymath}. De Bruijn showed that $H_t$ has only real zeros for
$t\ge\frac12$ \cite{debruijn}. Newman showed that there is a constant
$\Lambda\le\frac12$ such that $H_t$ has only real zeros exactly when
$t\ge\Lambda$ \cite{newman}, and the Riemann hypothesis is the statement
$\Lambda\le0$. Rodgers and Tao proved $\Lambda\ge0$ \cite{rodgerstao}. The
published upper bounds are $0.22$ \cite{polymath} and $0.2$
\cite{platttrudgian}. Dobner extended the construction to the extended
Selberg class, in particular to $\xi_\chi$ for primitive Dirichlet
characters, and proved $\Lambda_\chi\ge0$ \cite{dobner}.

A recent line of work follows the highest zero directly. If $x+ih$ is a
highest simple zero of $H_t$ and $Q=h^2$, the zero dynamics give
\begin{equation}\label{eq:qprime}
Q'\le-2-\frac{4Q}{p}\,S(x,p,t)+\frac{8Q}{p^2-Q}\qquad(p^2\ge5Q),\qquad
S(x,p,t)=-\im\frac{H_t'}{H_t}(x+ip)
\end{equation}
(\cite[Section~4]{hb}; Proposition~\ref{prop:paired} below). Hence a lower
bound for $S$ at one height $p$ and every relevant $x$ drives $Q$ down.
Hakimoglu-Brown \cite{hb} combines \eqref{eq:qprime} with three further
inputs: Polymath's finite-region computation, an approximation of
Riemann--Siegel type on the band $3\le\im z\le6$, and an upper bound for
$\partial_pS$ (his Proposition~11), which enters through a strengthened
kernel inequality (his Lemma~10 and Theorem~12). He obtains
$\Lambda\le0.13963$ under the hypothesis that $\zeta$ has no zero with
$\re s>\frac78$.

Every comparison of this kind needs a floor for $S(x,p,t)$ along the flow,
uniform over the real parts $x$ where the highest zero can lie. At $t=0$
the floor is exact. There $S(x,p,0)$ is half the
real part of $F'/F$ at $\re s=\frac{1+p}2>1$, which the Euler product
bounds below (Lemma~\ref{lem:floor0}). This paper is about keeping that
floor for $t>0$. In the language of the project it belongs to, it supplies
the exterior input of a strip comparison: the statement called
Lemma~$A_p$ in \cite{p01}. It does so without passing through a
zero-free region near the critical line.

\subsection{Results}

Throughout, $F$ is $\xi$ or $\xi_\chi$ for a primitive character $\chi$ of
conductor $q$ and parity $\fa$. Let $\Gamma_F$ be its Gamma factor and
$\alpha=\Gamma_F'/\Gamma_F$. Fix $p>1$ and $0<t\le\frac12$, and put
$\tau=\sqrt{t/2}$ and $w=\frac{1+p}2+\frac i2x$.

\begin{itemize}
\item \emph{Gaussian form and saddle (Section~\ref{sec:saddle}).} We have
$S(x,p,t)=\frac12\re(\cG'/\cG)(w)$, where $\cG(s)=\E\,F(s+\tau N)$ for a
standard normal variable $N$. The fixed point $s_*=w+\frac t2\alpha(s_*)$
exists and is unique (for $\zeta$, in a box $D_x$ at heights $x\ge40$,
for $p\le x$). Moving the line of integration to $s_*$ gives an
exact formula for $(\cG'/\cG)(w)-(F'/F)(s_*)$
(Lemmas~\ref{lem:gauss}--\ref{lem:shift}).
\item \emph{Explicit remainder (Section~\ref{sec:remainder}).} We have
$\bigl|(\cG'/\cG)(w)-(F'/F)(s_*)\bigr|\le E$, with $E$ given explicitly
by majorants at real arguments (Proposition~\ref{prop:E}). Consequently
\begin{equation}\label{eq:floorintro}
S(x,p,t)\ \ge\ \tfrac14\ell_*-\tfrac12\kappa(\sigma_*)-\tfrac12E,\qquad
\sigma_*=\tfrac{1+p}2+\tfrac t4\ell_*,\quad\ell_*=2\re\alpha(s_*),
\end{equation}
with $\kappa=-\zeta'/\zeta$ (Theorem~\ref{thm:transport}). The time-$t$
floor is therefore the time-zero floor at a point moved to the right by
$t\ell_*/4$; this is the \emph{height-$p$ transport}.
\item \emph{Derivative bound (Section~\ref{sec:derivative}).} Cauchy's
estimate applied to the remainder bounds $\partial_pS$
(Corollary~\ref{cor:derivative}).
\item \emph{Numerical evaluation (Section~\ref{sec:numerics}).} Floating
evaluation of $E$ on a grid of 20 times $0.001\le t\le\frac12$ shows the
floor of Lemma~\ref{lem:floor0} persisting with no loss. This holds for
every primitive character of conductor $q\ge1.084\cdot10^5$ (even) or
$q\ge4.892\cdot10^4$ (odd) at $p=\sqrt5$, at every height. For $\zeta$ the
constant floor $\frac14\ell_X-\frac12\kappa(\sigma_0)$ persists at every
$x\ge X$, for the heights $X$ used by Polymath \cite{polymath} and by Platt
and Trudgian \cite{platttrudgian}. At $p=61/20$ the floating remainder is
smaller than that of \cite[Proposition~11]{hb} by factors of $1.7$ to $4.4$.
\item \emph{A check from the classical strip
(Section~\ref{sec:strip}).} Conditionally on the analytic lemmas of
\cite{hb}, Hakimoglu-Brown's Theorem~12 started from the classical strip
$M(0)\le1$ gives $\Lambda\le18891/100000$ (Proposition~\ref{prop:table}).
The 34-row table is verified in exact and directed-rounding arithmetic.
\end{itemize}

Table~\ref{tab:status} records the status of each statement.

\begin{table}[ht]
\centering\small
\begin{tabularx}{\textwidth}{>{\raggedright\arraybackslash}p{0.42\textwidth}X}
\toprule
Statement & Status\\
\midrule
Lemmas~\ref{lem:floor0}--\ref{lem:partial}, Propositions~\ref{prop:paired} and~\ref{prop:E}, Theorem~\ref{thm:transport}, Corollary~\ref{cor:derivative} & proved here (Lemma~\ref{lem:floor0} and Proposition~\ref{prop:paired} after \cite{p01,hb})\\
Section~\ref{sec:numerics}: thresholds, floors, comparison with \cite{hb} & numerical: floating point, monotonicity in $\ell_*$ sampled, grid times only\\
Proposition~\ref{prop:table}: $\Lambda\le0.18891$ & conditional on \cite[Lemmas~2, 3, 6, 10, Proposition~11, Theorem~12]{hb} and the further inputs listed in Proposition~\ref{prop:table}\\
Section~\ref{sec:planning} & planning: depends on unproved inputs named there\\
\bottomrule
\end{tabularx}
\caption{Status of the statements.}\label{tab:status}
\end{table}

\subsection{Relation to other work}

The exterior inputs of Polymath \cite{polymath} and Hakimoglu-Brown
\cite{hb} come from approximate functional equations: a heat-weighted
Dirichlet polynomial, its reflection and Riemann--Siegel remainders. The
transport identity replaces these at heights $p>1$ by an exact Gaussian
average with a saddle-point recentring. Its main term coincides with
Hakimoglu-Brown's. His heat weight $e^{t\log^2n/4}$ equals
$\E\,n^{-\tau N}$, and his point $w+\frac t2\alpha(w)$ is, up to complex
conjugation, the first iterate of the saddle equation
(Section~\ref{sec:compare}).

Starting a stepped comparison from the classical strip is not new. Gordon
\cite{gordon}, Raymond \cite{raymond} and Planat \cite{planat} start from
it and claim $\Lambda\le0.157656619$, $\Lambda<0.149396$ and, under stated
premises, $\Lambda\le0.148611789321711$. Gomila \cite{gomila} claims
$\Lambda\le0.1787854$ by Polymath's criterion at a larger height. All four
are unrefereed, and we did not read them in full
(Section~\ref{sec:literature}). Proposition~\ref{prop:table} is a cross-check:
it rests on published numerical inputs and on the analytic lemmas of
\cite{hb}. For Dirichlet $L$-functions we know of no unconditional upper
bound below $\frac12$ for $\Lambda_\chi$ \cite{p01}. There the transport identity
is the main new input, and Section~\ref{sec:planning} records what it
would give.

\section{Setting, the time-zero floor and the comparison}\label{sec:setting}

\subsection{Conventions}\label{sec:conventions}

We write $s=\frac12(1+iz)$ with $z=x+iy$. A point $s=\beta+iT$ corresponds
to $z=2T+i(1-2\beta)$, so the real part $x$ is twice the ordinate and the
line $\re s=\sigma$ is the line $\im z=1-2\sigma$.

\emph{The family.} Let $\chi$ be a primitive Dirichlet character modulo
$q\ge3$ with $\chi(-1)=(-1)^{\fa}$, $\fa\in\{0,1\}$. Put $\cL=L(\cdot,\chi)$,
$\Gamma_F(s)=\epsilon_\chi(q/\pi)^{s/2}\Gamma(\frac{s+\fa}2)$ and
$F=\xi_\chi=\Gamma_F\cL$, where the constant $\epsilon_\chi$ of modulus one
is chosen so that $F(s)=\overline{F(1-\bar s)}$ \cite{dobner}.

\emph{The zeta case.} Here $\cL=\zeta$,
$\Gamma_F(s)=\frac12s(s-1)\pi^{-s/2}\Gamma(\frac s2)$, $F=\xi=\Gamma_F\zeta$,
and we put $q=1$ and $\fa=0$.

In both cases $\alpha=\Gamma_F'/\Gamma_F$ is given by
\begin{equation}\label{eq:alpha}
\alpha(s)=\tfrac12\log\tfrac q\pi+\tfrac12\psi\bigl(\tfrac{s+\fa}2\bigr)\quad\text{(family)},\qquad
\alpha(s)=\tfrac1s+\tfrac1{s-1}-\tfrac12\log\pi+\tfrac12\psi\bigl(\tfrac s2\bigr)\quad\text{(zeta)},
\end{equation}
with $\psi=\Gamma'/\Gamma$. Let
$H_t(z)=\int_\R e^{tu^2}\Phi(u)e^{izu}\dd u$, where $\Phi=\Phi_\chi$ is
Dobner's kernel \cite{dobner} for the family and Polymath's kernel
\cite{polymath} for $\zeta$. Then $\partial_tH_t=-\partial_z^2H_t$, $H_t$
is real on $\R$, and $H_0(z)=c_HF(\frac12(1+iz))$ with a real constant
$c_H\ne0$. (Polymath's $H_t$, defined by $\int_0^\infty\cos(zu)\dd u$, is
half of ours; constant factors do not change $S$.) The kernel decays like
$\exp(-\frac\pi qe^{4\lvert u\rvert})$, so each $H_t$ is entire of order
at most one.

Let $M(t)=\sup\{(\im z)^2:H_t(z)=0\}$. Every zero of $F$ has
$0<\re s<1$, so $M(0)\le1$. De Bruijn's strip theorem \cite{debruijn}
(for the family see \cite{dobner}) gives $M(t)\le\max(M(t_0)-2(t-t_0),0)$
for $t\ge t_0$, and in particular $\Lambda\le t_0+\frac12M(t_0)$. We call
this the \emph{classical strip theorem}. It implies that every zero of $H_t$,
$t\ge0$, satisfies $\lvert\im z\rvert\le1$.

Fix $p>1$ and put
\[
\sigma_0=\tfrac12(1+p),\qquad w=\sigma_0+\tfrac i2x,\qquad\tau=\sqrt{t/2},
\]
so that $z=x+ip$ corresponds to $s=1-\bar w$. Let $\varphi_\tau$ be the
density of the centred normal law of variance $\tau^2$, and let $N$ be a
standard normal variable. For $\sigma>1$ let
\[
\kappa(\sigma)=-\frac{\zeta'}{\zeta}(\sigma)=\sum_{n\ge2}\Lambda_{\mathrm{vM}}(n)n^{-\sigma},\qquad
K_2=-\kappa',\qquad K_3=\kappa'',
\]
with $\Lambda_{\mathrm{vM}}$ the von Mangoldt function. These are positive,
$\kappa$ and $K_2$ decrease, and $\log\zeta$ and $\kappa$ are convex.

\subsection{The time-zero floor}

\begin{lemma}[time-zero floor; after {\cite[Lemma~3.4]{p01}}]\label{lem:floor0}
Let $p>1$. For the family, at every real $x$,
\begin{equation}\label{eq:Cq}
S(x,p,0)\ \ge\ C_{q,\fa}(p):=\tfrac14\Bigl(\log\tfrac q\pi+\psi\bigl(\tfrac{\sigma_0+\fa}2\bigr)\Bigr)-\tfrac12\kappa(\sigma_0).
\end{equation}
For $\zeta$, at every $x>0$,
$S(x,p,0)\ge\frac12\re\alpha(w)-\frac12\kappa(\sigma_0)$, and
$\frac12\re\alpha(w)\ge\frac14\log\frac x{4\pi}-\frac14c(\frac w2)$ with
$c$ as in Lemma~\ref{lem:stirling}.
\end{lemma}

\begin{proof}
By Lemma~\ref{lem:gauss} at $t=0$, $S(x,p,0)=\frac12\re(F'/F)(w)$, and
$F'/F=\alpha+\cL'/\cL$. Since
$\cL'/\cL(s)=-\sum\chi(n)\Lambda_{\mathrm{vM}}(n)n^{-s}$,
$\lvert(\cL'/\cL)(w)\rvert\le\kappa(\sigma_0)$. For $\re z>0$,
$\psi(z)=-\gamma_E+\sum_{n\ge0}\bigl(\frac1{n+1}-\frac1{n+z}\bigr)$ and
$\re\frac1{n+z}\le\frac1{n+\re z}$, so $\re\psi(z)\ge\psi(\re z)$. This
gives the family bound. For $\zeta$, $\re(\frac1w+\frac1{w-1})>0$ and, by
Lemma~\ref{lem:stirling} with $z=w/2$,
$\re\psi(\frac w2)\ge\log\frac{\lvert w\rvert}2-c(\frac w2)$. Hence
$\re\alpha(w)\ge\frac12\log\frac{\lvert w\rvert}{2\pi}-\frac12c(\frac w2)$,
and $\lvert w\rvert\ge\frac x2$.
\end{proof}

\subsection{The paired-zero inequality}

For $p,h>0$, $d\in\R$ and $0\le v$, let
\[
K_p(d,v)=\frac{p-v}{d^2+(p-v)^2}+\frac{p+v}{d^2+(p+v)^2}
=\frac{2p\,(d^2+p^2-v^2)}{[d^2+(p-v)^2]\,[d^2+(p+v)^2]}.
\]
For a zero $z_k=x_k+iy_k$ with $y_k\ge0$ write $d_k=x-x_k$, $v_k=y_k$,
and give it the weight $\varpi_k=1$ if $y_k>0$ and $\varpi_k=\frac12$ if
$y_k=0$. Zeros are counted with multiplicity.

\begin{proposition}[paired-zero inequality; after {\cite[Section~4]{hb}} and {\cite[Proposition~3.3]{p01}}]\label{prop:paired}
Let $p>1$ and $t>0$.
\begin{enumerate}
\item[(i)] $S(x,p,t)=\sum_{y_k\ge0}\varpi_kK_p(d_k,v_k)$, summed over the
zeros of $H_t$; the series converges absolutely.
\item[(ii)] Let $z_j=x+ih$, $h>0$, be a simple zero of $H_t$. Its
trajectory is differentiable in $t$, and
$h'=-\frac1h-2\sum'\varpi_kK_h(d_k,v_k)$, where $\sum'$ omits $z_j$ and
$\bar z_j$.
\item[(iii)] If moreover no zero of $H_t$ has imaginary part larger than
$h$, and $p^2\ge5h^2$, then $Q=h^2$ satisfies \eqref{eq:qprime}.
\end{enumerate}
\end{proposition}

\begin{proof}
Since $H_t$ has order at most one and is real on $\R$, its zeros are
closed under conjugation and
$H_t(z)=z^me^{a+bz}\prod_k(1-z/z_k)e^{z/z_k}$ with $b$ real. Then
$(H_t'/H_t)(z)=\frac mz+b+\sum_k\bigl(\frac1{z-z_k}+\frac1{z_k}\bigr)$,
absolutely convergent. Group each nonreal $z_k$ with $\bar z_k$. The
terms $\frac1{z_k}+\frac1{\bar z_k}$ are real, and
$\im\bigl(\frac1{z-z_k}+\frac1{z-\bar z_k}\bigr)=-K_p(d_k,v_k)$ at
$z=x+ip$. A real zero contributes $\im\frac1{z-x_k}=-\frac12K_p(d_k,0)$.
This proves (i). For (ii), the implicit function theorem gives a
differentiable trajectory with
$z_j'=-\partial_tH_t/\partial_zH_t=\partial_z^2H_t/\partial_zH_t$ at $z_j$.
Writing $H_t=(z-z_j)g$, this equals
$2(g'/g)(z_j)=2\bigl[\frac m{z_j}+b+\frac1{z_j}+\sum_{k\ne j}\bigl(\frac1{z_j-z_k}+\frac1{z_k}\bigr)\bigr]$.
In the imaginary part the terms $\frac1{z_j}$ and $\frac1{\bar z_j}$
combine to a real number, the conjugate zero contributes
$\im\frac2{z_j-\bar z_j}=-\frac1h$, each other pair contributes
$-2K_h(d_k,v_k)$, and each real zero $-K_h(d_k,0)$, as in (i).

For (iii) we use the kernel inequality $K_h(d,v)\ge\frac hpK_p(d,v)$
for $0\le v\le h$, $(d,v)\ne(0,h)$ and $p^2\ge5h^2$. With $D=d^2$,
$A=h^2$, $B=p^2$, $r=v^2\le A$, $w_A=D+A-r$ and $w_B=D+B-r$, the
denominators of $K_h$ and $K_p$ are $w_A^2+4Dr$ and $w_B^2+4Dr$, and
\[
\frac{K_h}h-\frac{K_p}p=\frac{2(B-A)\bigl(w_Aw_B-4Dr\bigr)}{(w_A^2+4Dr)(w_B^2+4Dr)},\qquad
w_Aw_B-4Dr=D^2+(A+B-6r)D+(A-r)(B-r),
\]
which is nonnegative because $A+B-6r\ge B-5A\ge0$. Every other zero has
$v_k\le h$, so by (i) and (ii)
\[
Q'=2hh'=-2-4h{\sum}'\varpi_kK_h\le-2-\frac{4Q}p\Bigl(S(x,p,t)-K_p(0,h)\Bigr),
\qquad K_p(0,h)=\frac{2p}{p^2-Q},
\]
which is \eqref{eq:qprime}.
\end{proof}

Hakimoglu-Brown notes that the kernel inequality also appears in Gordon's
preprint \cite[Section~5]{gordon}. Lemma~\ref{lem:hb10} below strengthens
it by a term in $\partial_pK_p$, which Theorem~\ref{thm:hb12} pays for with
an upper bound for $\partial_pS$.

\subsection{The exterior input}\label{sec:exterior}

A comparison built on \eqref{eq:qprime} needs, at one fixed height $p$,
\begin{equation}\label{eq:Ep}
\textbf{(E$_p$)}\qquad S(x,p,t)\ \ge\ C(t)\qquad(0\le t\le t_*,\ x\text{ in the region where the highest zero can lie}).
\end{equation}
For the family at large conductor the region is the whole line. For $\zeta$
and for small conductors it is $\lvert x\rvert\ge X$, and the rest is
covered by a finite verification and a barrier, as in
\cite[Lemma~2]{hb}. In \cite{p01}, \eqref{eq:Ep} is called Lemma~$A_p$. The
comparison also needs uniform confinement of the nonreal zeros at positive
times (\cite{p01}, item H12) and a first-contact argument; for $\zeta$,
\cite{hb} supplies these through \cite[Theorem~1.5]{polymath} and its
Lemma~2. The theorems of this paper concern \eqref{eq:Ep} only.
Statement~\eqref{eq:Ep} lives at height $p>1$, where the Dirichlet series
converges absolutely: it bounds a Poisson-smoothed density of zeros from
below, and uses no zero-free region and no derivative of $H_t$ beyond the
first.

\section{The transport identity}\label{sec:saddle}

\begin{lemma}[Gaussian form]\label{lem:gauss}
For $t\ge0$ and $z\in\C$, $H_t(z)=\E\,H_0(z+i\sqrt{2t}\,N)$. With
$\cG(s)=\E\,F(s+\tau N)$ one has $H_t(z)=c_H\,\cG(\frac12(1+iz))$,
$\cG(s)=\overline{\cG(1-\bar s)}$, and, for $p>1$,
\[
S(x,p,t)=\tfrac12\re\frac{\cG'}{\cG}(w).
\]
\end{lemma}

\begin{proof}
By Fubini, which the decay of $\Phi$ permits,
$\E\,H_0(z+i\beta N)=\int\Phi(u)e^{izu}\E\,e^{-\beta Nu}\dd u=\int e^{\beta^2u^2/2}\Phi(u)e^{izu}\dd u$,
and $\beta=\sqrt{2t}$ gives $H_t$. Substituting $H_0(z)=c_HF(\frac12(1+iz))$
and using $N\stackrel d=-N$ gives $H_t(z)=c_H\cG(\frac12(1+iz))$ with
$\tau=\beta/2$, and the functional equation of $F$ passes to $\cG$. Then
$H_t'(z)=\frac i2c_H\cG'(\frac12(1+iz))$. At $z=x+ip$,
$\frac12(1+iz)=1-\bar w$, and differentiating
$\cG(s)=\overline{\cG(1-\bar s)}$ gives
$(\cG'/\cG)(1-\bar w)=-\overline{(\cG'/\cG)(w)}$. So
$S=-\im\bigl(\frac i2(\cG'/\cG)(1-\bar w)\bigr)=\frac12\re(\cG'/\cG)(w)$.
Here $\cG(w)\ne0$, because every zero of $H_t$ has $\lvert\im z\rvert\le1<p$.
\end{proof}

\begin{lemma}[the Gamma saddle]\label{lem:saddle}
\emph{Family.} Let $p>1$, $0<t\le\frac12$ and
$\ell_0=\log\frac q\pi+\psi(\frac{\sigma_0+\fa}2)>0$. The map
$\mu\mapsto\tau^2\alpha(w+\mu)$ sends the closed half-plane
$\{\re\mu\ge0\}$ into itself and is a contraction there, so it has a unique
fixed point $\mu_*$. Put $s_*=w+\mu_*$, $\ell_*=2\re\alpha(s_*)$ and
$\sigma_*=\re s_*$. Then
\[
\sigma_*=\sigma_0+\tfrac12\tau^2\ell_*=\sigma_0+\tfrac t4\ell_*,\qquad
\ell_*\ge\ell_*^{(0)}\ge\ell_0,
\]
where $\ell_*^{(0)}=2\alpha(\sigma^{(0)})$ and $\sigma^{(0)}$ is the unique
real solution of $\sigma=\sigma_0+\tau^2\alpha(\sigma)$, the saddle at $x=0$.

\emph{Zeta.} For $x\ge40$ and $1<p\le x$ the same map has a unique fixed point
in $D_x=\{\mu:0\le\re\mu\le R_x,\ \lvert\im\mu\rvert\le\frac12\}$,
$R_x=\frac18\log x+1$. There $\im s_*\ge\frac x2-\frac12$ and
$\ell_*\ge\log\bigl((\frac x2-\frac12)/2\pi\bigr)-2c(\frac12s_*)$, with $c$
as in Lemma~\ref{lem:stirling}. ($S$ is even in $x$ for $\zeta$.)
\end{lemma}

\begin{proof}
Family. For $\re\mu\ge0$, $\re(w+\mu)\ge\sigma_0$. As in the proof of
Lemma~\ref{lem:floor0}, $\re\psi(z)\ge\psi(\re z)$, and termwise from
$\psi'(z)=\sum_{n\ge0}(n+z)^{-2}$, $\lvert\psi'(z)\rvert\le\psi'(\re z)$.
So $\re\tau^2\alpha(w+\mu)\ge\frac12\tau^2\ell_0>0$ and
$\bigl\lvert\frac{\dd}{\dd\mu}\tau^2\alpha(w+\mu)\bigr\rvert\le\frac14\tau^2\psi'(\frac{\sigma_0+\fa}2)\le\frac1{16}\psi'(\frac12)=\frac{\pi^2}{32}<0.31$
(below $0.15$ when $p\ge\sqrt5$). Banach's theorem gives $\mu_*$, and
$\re\mu_*=\tau^2\re\alpha(s_*)=\frac12\tau^2\ell_*$. Next,
$\ell_*\ge2\alpha(\sigma_*)$ and $\sigma_*\ge\sigma_0+\tau^2\alpha(\sigma_*)$.
The function $f(\sigma)=\sigma-\sigma_0-\tau^2\alpha(\sigma)$ is strictly
increasing ($\tau^2\alpha'<1$) with $f(\sigma^{(0)})=0\le f(\sigma_*)$, so
$\sigma_*\ge\sigma^{(0)}\ge\sigma_0$, and $\alpha$ increases on the real
axis.

Zeta. For $\mu\in D_x$, $s=w+\mu$ has $\re s\in[\sigma_0,\sigma_0+R_x]$ and
$\im s\ge\frac x2-\frac12\ge19.5$. There $\lvert\alpha'(s)\rvert\le2/\im s\le0.11$
by Lemma~\ref{lem:gamma}, so the map is a contraction. With $z=s/2$,
$\im\psi(z)=\sum_{n\ge0}\im z/\lvert z+n\rvert^2\in(0,\frac\pi2+\frac1{2\re z}]$
and $\lvert\im(\frac1s+\frac1{s-1})\rvert\le2/\im s$, so
$\lvert\im\tau^2\alpha(s)\rvert\le\frac14(\frac\pi4+\frac12+0.11)<\frac12$.
Lemma~\ref{lem:stirling} and $\re(\frac1s+\frac1{s-1})>0$ give
$\frac12\log\frac{\lvert s\rvert}{2\pi}-c(z)\le\re\alpha(s)\le\frac12\log\frac{\lvert s\rvert}{2\pi}+\frac1{12\re z\lvert z\rvert}+\frac2{\im s}$,
the upper bound from $\re\psi(z)\le\log\lvert z\rvert+\lvert I(z)\rvert$ in
the notation of that lemma. The lower bound is positive. Since $p\le x$,
$\lvert s\rvert\le\re s+\im s\le x+2+\frac18\log x\le2x$ on $D_x$, so
$\tau^2\re\alpha(s)\le\frac18\log\frac x\pi+0.03<R_x$. So $D_x$ is
mapped into itself. The bound on $\ell_*$ is the lower bound at $s_*$.
\end{proof}

\begin{lemma}[the shifted representation]\label{lem:shift}
Let $\Theta(v)=F(s_*+v)e^{-\alpha(s_*)v}/F(s_*)$, an entire function of $v$
of order one with $\Theta(0)=1$, and $c_1=(\cL'/\cL)(s_*)$. (Since
$\re s_*>1$, $F(s_*)\ne0$.) Then
\[
\frac{\cG'}{\cG}(w)-\frac{F'}{F}(s_*)=\frac{\int_\R(\Theta'-c_1\Theta)\varphi_\tau}{\int_\R\Theta\varphi_\tau},\qquad
\int_\R(\Theta'-c_1\Theta)\varphi_\tau=\int_\R\Bigl(\frac v{\tau^2}-c_1\Bigr)\Theta\varphi_\tau .
\]
\end{lemma}

\begin{proof}
The function $h\mapsto F^{(k)}(w+h)\varphi_\tau(h)$ is entire, and on
$\re h=\pm R$, between $0$ and $\im\mu_*$, it is
$O(e^{R^{1+\epsilon}-R^2/2\tau^2})$. So the line of integration in
$\cG^{(k)}(w)=\int F^{(k)}(w+h)\varphi_\tau(h)\dd h$ moves to
$\mu_*+\R$. Since $\mu_*/\tau^2=\alpha(s_*)$,
$\varphi_\tau(\mu_*+v)=c_0\varphi_\tau(v)e^{-\alpha(s_*)v}$ with
$c_0=e^{-\mu_*^2/2\tau^2}$. Hence $\cG(w)=c_0F(s_*)\int\Theta\varphi_\tau$
and $\cG'(w)=c_0F(s_*)\int(\Theta'+\alpha(s_*)\Theta)\varphi_\tau$, because
$F'(s_*+v)e^{-\alpha(s_*)v}/F(s_*)=\Theta'(v)+\alpha(s_*)\Theta(v)$.
Subtract $(F'/F)(s_*)=\alpha(s_*)+c_1$. The second identity is Gaussian
integration by parts, $\varphi_\tau'=-(v/\tau^2)\varphi_\tau$, valid because
$\Theta$ has order one.
\end{proof}

The question is therefore how far the complex weight $\Theta\varphi_\tau$
is from the pure tilt $e^{c_1v}\varphi_\tau$, for which the right side
vanishes. Where $\re(s_*+v)>1$ we have $\Theta=e^{c_1v+\omega(v)}$ with
\[
\omega(v)=\bigl[\log\cL(s_*+v)-\log\cL(s_*)-c_1v\bigr]+\bigl[\log\Gamma_F(s_*+v)-\log\Gamma_F(s_*)-\alpha(s_*)v\bigr],
\]
and $\Theta'-c_1\Theta=\Delta\Theta$ with
$\Delta(v)=(F'/F)(s_*+v)-(F'/F)(s_*)$.

\section{An explicit remainder}\label{sec:remainder}

For a function $f$ of a real variable write
$\Rem{f}_v=f(\sigma_*+v)-f(\sigma_*)-vf'(\sigma_*)$ for its Taylor
remainder at $\sigma_*$.

\begin{lemma}[Euler-product majorants]\label{lem:euler}
For $\re s=\sigma>1$ and $j\ge0$,
$\lvert(\cL'/\cL)^{(j)}(s)\rvert\le(-1)^j\kappa^{(j)}(\sigma)$,
$\lvert\log\cL(s)\rvert\le\log\zeta(\sigma)$ and
$\zeta(2\sigma)/\zeta(\sigma)\le\lvert\cL(s)\rvert\le\zeta(\sigma)$. Hence,
for real $v$ with $\sigma_*+\min(v,0)>1$,
\begin{align*}
\lvert(\cL'/\cL)(s_*+v)-c_1\rvert&\le\lvert\kappa(\sigma_*+v)-\kappa(\sigma_*)\rvert,\\
\lvert(\cL'/\cL)(s_*+v)-c_1-v(\cL'/\cL)'(s_*)\rvert&\le\Rem{\kappa}_v,\\
\lvert\log\cL(s_*+v)-\log\cL(s_*)\rvert&\le\lvert\log\zeta(\sigma_*+v)-\log\zeta(\sigma_*)\rvert,\\
\lvert\log\cL(s_*+v)-\log\cL(s_*)-c_1v\rvert&\le\Rem{\log\zeta}_v .
\end{align*}
\end{lemma}

\begin{proof}
$\cL'/\cL=-\sum\chi(n)\Lambda_{\mathrm{vM}}(n)n^{-s}$ and
$\log\cL=\sum\chi(n)\Lambda_{\mathrm{vM}}(n)(\log n)^{-1}n^{-s}$ with
$\lvert\chi(n)\rvert\le1$; the bounds on $\lvert\cL\rvert$ follow from the
Euler product. Each difference is an integral along the horizontal segment
from $s_*$ to $s_*+v$ of a derivative bounded by a majorant of fixed sign,
whose integral is the corresponding difference or Taylor remainder of
$\kappa$ or $\log\zeta$; both are convex.
\end{proof}

\begin{lemma}[Gamma majorants]\label{lem:gamma}
\emph{Family.} For $\re z>0$, $\lvert\psi'(z)\rvert\le\psi'(\re z)$ and
$\lvert\psi''(z)\rvert\le-\psi''(\re z)$. Put
$\gamma_\Gamma(\sigma)=-\frac12\psi(\frac{\sigma+\fa}2)$ and
$\lambda_\Gamma(\sigma)=\log\Gamma(\frac{\sigma+\fa}2)$. For real $v$ with
$\sigma_*+\min(v,0)>-\fa$,
\begin{gather*}
\lvert\alpha(s_*+v)-\alpha(s_*)\rvert\le\lvert\gamma_\Gamma(\sigma_*+v)-\gamma_\Gamma(\sigma_*)\rvert,\qquad
\lvert\alpha(s_*+v)-\alpha(s_*)-v\alpha'(s_*)\rvert\le\Rem{\gamma_\Gamma}_v,\\
\lvert\log\Gamma_F(s_*+v)-\log\Gamma_F(s_*)-\alpha(s_*)v\rvert\le\Rem{\lambda_\Gamma}_v=:\vartheta(v).
\end{gather*}
\emph{Zeta.} If $\lvert\im s\rvert\ge T\ge10$, at any real part,
$\lvert\alpha'(s)\rvert\le2/T$ and $\lvert\alpha''(s)\rvert\le3/T^2$. Put
$\vartheta(v)=v^2/T$.
\end{lemma}

\begin{proof}
The family bounds are integrals along the segment, as in
Lemma~\ref{lem:euler}, since $\lvert\alpha'\rvert\le-\gamma_\Gamma'$,
$\lvert\alpha''\rvert\le\gamma_\Gamma''$ and $\lambda_\Gamma''=-\gamma_\Gamma'$.
For $\zeta$, with $z=s/2=a+ib$, the summands of
$\sum_{n\ge0}\lvert z+n\rvert^{-2}$ and $\sum_{n\ge0}\lvert z+n\rvert^{-3}$ are
unimodal in $n$, so each sum is at most the integral plus the largest term:
$\sum_{n\ge0}\lvert z+n\rvert^{-2}\le b^{-2}+\pi/\lvert b\rvert$ and
$\sum_{n\ge0}\lvert z+n\rvert^{-3}\le\lvert b\rvert^{-3}+2b^{-2}$, with
$\lvert b\rvert\ge T/2$. Hence
$\frac14\lvert\psi'(z)\rvert\le\frac\pi{2T}+\frac1{T^2}$ and
$\frac18\lvert\psi''(z)\rvert\le\frac2{T^3}+\frac2{T^2}$. Adding
$\lvert s\rvert^{-2}+\lvert s-1\rvert^{-2}\le2T^{-2}$ and
$2\lvert s\rvert^{-3}+2\lvert s-1\rvert^{-3}\le4T^{-3}$ gives
$\lvert\alpha'\rvert\le\frac\pi{2T}+\frac3{T^2}\le\frac2T$ and
$\lvert\alpha''\rvert\le\frac2{T^2}+\frac6{T^3}\le\frac3{T^2}$ for $T\ge10$.
\end{proof}

\begin{lemma}[a Stirling inequality]\label{lem:stirling}
For $\re z>0$,
\[
\log\lvert z\rvert\ \le\ \re\psi(z)+c(z),\qquad c(z)=\frac{\re z}{2\lvert z\rvert^2}+\frac1{12\re z\,\lvert z\rvert}.
\]
\end{lemma}

\begin{proof}
Binet's formula \cite[Section~5.9(ii)]{dlmf} gives
$\psi(z)=\log z-\frac1{2z}-I(z)$ with $I(z)=\int_0^\infty\beta(u)e^{-zu}\dd u$
and $\beta(u)=\frac1{e^u-1}-\frac1u+\frac12$. Since $\beta(0)=0$,
$I(z)=z^{-1}\int_0^\infty\beta'(u)e^{-zu}\dd u$, and
$\beta'(u)=\frac14f(u/2)$ with $f(v)=v^{-2}-\sinh^{-2}v$. Here $f>0$
because $\sinh v>v$, and $f\le\frac13$: for $v\ge\sqrt3$ because
$v^{-2}\le\frac13$, and for $v<\sqrt3$ because
$\sinh^2v\,(3-v^2)=\sum_{n\ge1}(3a_n-a_{n-1})v^{2n}$ with
$a_n=2^{2n-1}/(2n)!$ and $a_0=0$, whose coefficients are $3$, $0$ and then
$a_{n-1}\bigl(\frac{12}{2n(2n-1)}-1\bigr)<0$, so that
$\sinh^2v\,(3-v^2)\le3v^2$. Hence $0\le\beta'\le\frac1{12}$,
$\lvert I(z)\rvert\le1/(12\re z\,\lvert z\rvert)$, and taking real parts
gives the claim.
\end{proof}

For the family put $c_*=\frac1{2x_*}+\frac1{12x_*^2}$ with
$x_*=\frac12(\sigma_*+\fa)$. It bounds $c(z_*)$, $z_*=\frac12(s_*+\fa)$, at
every height. Since $\ell_*=\log\frac q\pi+\re\psi(z_*)$,
Lemma~\ref{lem:stirling} gives $\log\lvert z_*\rvert\le\ell_*-\log\frac q\pi+c_*$,
so for $0<u\le\sigma_*$
\begin{equation}\label{eq:Qstar}
q\,\lvert u+i\im s_*\rvert\ \le\ 2q\lvert z_*\rvert\ \le\ \mathcal Q_*:=2\pi e^{\ell_*+c_*}.
\end{equation}
For $\zeta$ the same $\mathcal Q_*$, with $c_*=c(\frac12s_*)$, bounds
$\lvert u+i\im s_*\rvert$.

\begin{lemma}[partial summation in the strip]\label{lem:partial}
For $\sigma>0$, $s=\sigma+iT$ and every integer $Y\ge1$,
\begin{align*}
\lvert L(s,\chi)\rvert&\le1+\int_1^Yu^{-\sigma}\dd u+\frac{q\lvert s\rvert Y^{-\sigma}}{2\sigma}
\qquad(\chi\text{ primitive modulo }q\ge3),\\
\lvert\zeta(s)\rvert&\le1+\int_1^Yu^{-\sigma}\dd u+\frac{Y^{1-\sigma}}{\lvert s-1\rvert}+\frac12Y^{-\sigma}+\frac{\lvert s\rvert Y^{-\sigma}}{2\sigma}.
\end{align*}
\end{lemma}

\begin{proof}
$L(s,\chi)=\sum_{n\le Y}\chi(n)n^{-s}+s\int_Y^\infty(A(u)-A(Y))u^{-s-1}\dd u$
with $A(u)=\sum_{n\le u}\chi(n)$. A sum of $\chi$ over consecutive integers
equals a sum over $r<q$ of them, which is also minus the sum over the other
$q-r$ of a period, so $\lvert A(u)-A(Y)\rvert\le q/2$. For $\zeta$,
$\zeta(s)=\sum_{n\le Y}n^{-s}+\frac{Y^{1-s}}{s-1}-\frac12Y^{-s}-s\int_Y^\infty(\{u\}-\frac12)u^{-s-1}\dd u$.
Both identities hold for $\sigma>1$ and extend to $\sigma>0$.
\end{proof}

In the evaluation we take $Y$ near $\mathcal Q_*/2$, and for $\zeta$ we use
$\lvert s-1\rvert\ge\lvert\im s\rvert\ge T$. Lemma~\ref{lem:partial} is
weaker than the convexity bound by about $(q\lvert s\rvert)^{(1-\sigma)/2}$.
An explicit convexity bound would improve the small-conductor thresholds
of Section~\ref{sec:numerics}; we do not use one, so no unread import
enters.

\begin{proposition}[the remainder]\label{prop:E}
Assume the hypotheses of Lemma~\ref{lem:saddle}, and for $\zeta$ let
$10\le T\le\im s_*$ (for instance $T=\frac x2-\frac12$). Let
$t\in(0,\frac12]$ and $V\in(0,\sigma_*-1)$. Write
\[
k_*=K_2(\sigma_*)+A_1(\sigma_*),\quad A_1(\sigma)=\tfrac14\psi'\bigl(\tfrac{\sigma+\fa}2\bigr)\qquad(\text{zeta: }k_*=K_2(\sigma_*)+2/T),
\]
$\kappa_*=\kappa(\sigma_*)$, and $\mathfrak M=\kappa+\gamma_\Gamma$ (zeta:
$\mathfrak M=\kappa$). For $v\ge-V$ define
\begin{gather*}
r(v)=\Rem{\mathfrak M}_v\quad\bigl(\text{zeta: }\Rem{\kappa}_v+\tfrac32v^2/T^2\bigr),\qquad
\omega_m(v)=\Rem{\log\zeta}_v+\vartheta(v),\\
g(v)=\exp\bigl(\lvert\log\zeta(\sigma_*+v)-\log\zeta(\sigma_*)\rvert+\vartheta(v)\bigr),
\end{gather*}
and for $v<-V$, with $\sigma=\sigma_*+v$ and
$z_\zeta=\zeta(\sigma_*)/\zeta(2\sigma_*)$,
\[
g_L(v)=\begin{cases}z_\zeta\,\mathcal Z(\sigma)\,e^{\vartheta(v)},&\frac12\le\sigma<\sigma_*-V,\\
z_\zeta\,\mathcal Z(1-\sigma)\exp\bigl(\frac12\ell_*(1-2\sigma)+\vartheta(1-\sigma-\sigma_*)\bigr),&\sigma<\frac12.\end{cases}
\]
Here $\mathcal Z(u)=\zeta(u)$ for $u\ge\sigma_*$; for $1<u<\sigma_*$ it is
the smaller of $\zeta(u)$ and the bound of Lemma~\ref{lem:partial} at real
part $u$, for a fixed integer $Y\ge1$, with $q\lvert s\rvert$ (zeta:
$\lvert s\rvert$) replaced by $\mathcal Q_*$ and, for $\zeta$,
$\lvert s-1\rvert$ by $T$; and for $u\le1$ it is that bound alone. Put
\begin{gather*}
M_\rho=\int_{-V}^\infty r\,g\,\varphi_\tau,\qquad B_V=g(-V)\varphi_\tau(V),\qquad
P_1=\int_{-V}^\infty e^{\kappa_*\lvert v\rvert}\bigl(e^{\omega_m}-1\bigr)\varphi_\tau,\\
D_L=\int_{-\infty}^{-V}g_L\varphi_\tau,\qquad
N_L=\int_{-\infty}^{-V}g_L\Bigl(\frac{\lvert v\rvert}{\tau^2}+\kappa_*\Bigr)\varphi_\tau,\qquad
P_2=\int_{-\infty}^{-V}\bigl(g_L+e^{\kappa_*\lvert v\rvert}\bigr)\varphi_\tau,
\end{gather*}
and $d_0=e^{-\tau^2\kappa_*^2/2}-P_1-P_2$. If $\tau^2k_*<1$ and $d_0>0$,
then
\begin{equation}\label{eq:E}
\Bigl\lvert\frac{\cG'}{\cG}(w)-\frac{F'}{F}(s_*)\Bigr\rvert\ \le\ E:=
\frac{\tau^2k_*\kappa_*}{1-\tau^2k_*}+\frac{M_\rho+B_V+\tau^2k_*\kappa_*D_L+(1+\tau^2k_*)N_L}{(1-\tau^2k_*)\,d_0}.
\end{equation}
\end{proposition}

\begin{proof}
Write $\mathcal N$ and $\mathcal D$ for the numerator and denominator of
Lemma~\ref{lem:shift}. On $[-V,\infty)$, Lemmas~\ref{lem:euler}
and~\ref{lem:gamma} give $\lvert\Theta\rvert\le g$ and
$\lvert\omega\rvert\le\omega_m$. With $\Delta_1=\Delta'(0)=(F'/F)'(s_*)$
they also give $\lvert\Delta_1\rvert\le k_*$ and $\lvert\rho\rvert\le r$ for
$\rho(v)=\Delta(v)-\Delta_1v$.

For $v<-V$ we have $\lvert\Theta(v)\rvert\le g_L(v)$. When
$\sigma\ge\frac12$,
$\lvert\Theta\rvert=\lvert\cL(\sigma+i\im s_*)/\cL(s_*)\rvert\exp\re\bigl[\log\Gamma_F(s_*+v)-\log\Gamma_F(s_*)-\alpha(s_*)v\bigr]$,
and Lemmas~\ref{lem:euler}--\ref{lem:partial} and \eqref{eq:Qstar} apply,
with $\lvert\cL(s_*)\rvert\ge\zeta(2\sigma_*)/\zeta(\sigma_*)$. When
$\sigma<\frac12$, $\lvert F(\sigma+iT')\rvert=\lvert F(1-\sigma+iT')\rvert$
at $T'=\im s_*$. The Gamma ratio from $s_*$ to $1-\sigma+i\im s_*$
contributes $\exp\bigl(\frac12\ell_*(1-\sigma-\sigma_*)+\vartheta(1-\sigma-\sigma_*)\bigr)$,
which with $\lvert e^{-\alpha(s_*)v}\rvert=e^{\ell_*(\sigma_*-\sigma)/2}$
gives the stated exponent.

Gaussian integration by parts on $[-V,\infty)$,
$\int_{-V}^\infty v\varphi_\tau\Theta=\tau^2\varphi_\tau(V)\Theta(-V)+\tau^2\int_{-V}^\infty\varphi_\tau\Theta'$,
turns Lemma~\ref{lem:shift} into
$\mathcal N=\varphi_\tau(V)\Theta(-V)+\mathcal J+\mathcal N_L$ with
$\mathcal J=\int_{-V}^\infty\Delta\Theta\varphi_\tau$ and
$\mathcal N_L=\int_{-\infty}^{-V}(v/\tau^2-c_1)\Theta\varphi_\tau$. Applied
once more to $\Delta_1\int_{-V}^\infty v\Theta\varphi_\tau$, it gives
\[
(1-\tau^2\Delta_1)\mathcal J=\tau^2\Delta_1\bigl(\varphi_\tau(V)\Theta(-V)+c_1\mathcal D_E\bigr)+\int_{-V}^\infty\rho\Theta\varphi_\tau,\qquad
\mathcal D_E=\int_{-V}^\infty\Theta\varphi_\tau .
\]
Hence, with $\mathcal D_L=\int_{-\infty}^{-V}\Theta\varphi_\tau$ and
$\mathcal D=\mathcal D_E+\mathcal D_L$,
\[
\mathcal N-\frac{\tau^2\Delta_1c_1}{1-\tau^2\Delta_1}\mathcal D=
\frac{\varphi_\tau(V)\Theta(-V)+\int_{-V}^\infty\rho\Theta\varphi_\tau-\tau^2\Delta_1c_1\mathcal D_L+(1-\tau^2\Delta_1)\mathcal N_L}{1-\tau^2\Delta_1},
\]
whose modulus is at most $d_0$ times the second fraction of $E$, because
$\lvert c_1\rvert\le\kappa_*$, $\lvert\mathcal D_L\rvert\le D_L$ and
$\lvert\mathcal N_L\rvert\le N_L$. Finally
\[
\mathcal D=e^{\tau^2c_1^2/2}+\int_{-V}^\infty e^{c_1v}\bigl(e^{\omega}-1\bigr)\varphi_\tau+\int_{-\infty}^{-V}\bigl(\Theta-e^{c_1v}\bigr)\varphi_\tau
\]
and $\lvert e^{\tau^2c_1^2/2}\rvert\ge e^{-\tau^2\kappa_*^2/2}$, so
$\lvert\mathcal D\rvert\ge d_0$. Divide.
\end{proof}

\begin{theorem}[transport law and floor]\label{thm:transport}
Under the hypotheses of Proposition~\ref{prop:E},
\[
\Bigl\lvert S(x,p,t)-\tfrac12\re\tfrac{F'}F(s_*)\Bigr\rvert\le\tfrac12E,\qquad
S(x,p,t)\ \ge\ \tfrac14\ell_*-\tfrac12\kappa(\sigma_*)-\tfrac12E .
\]
For the family every quantity in $E$ is a function of
$(\ell_*,\tau,p,\fa,V)$, so the bound at $x$ is $\mathfrak F(\ell_*(x))$ with
\[
\mathfrak F(\ell)=\tfrac14\ell-\tfrac12\kappa\bigl(\sigma_0+\tfrac12\tau^2\ell\bigr)-\tfrac12E(\ell).
\]
For $\zeta$, $E$ depends on $\ell_*$ and, through $T$ and through $c_*$ in
$\mathcal Q_*$, on the height; $c_*$ may be bounded by its value at
$\re z=\frac12\sigma_*$, $\lvert z\rvert=\frac12(\sigma_*^2+T^2)^{1/2}$.
\end{theorem}

\begin{proof}
Lemma~\ref{lem:gauss}, Proposition~\ref{prop:E}, and
$\frac12\re(F'/F)(s_*)=\frac14\ell_*+\frac12\re c_1\ge\frac14\ell_*-\frac12\kappa_*$.
In the family case $\mathcal Q_*$, $\mathcal Z$, $\vartheta$, $r$, $g$ and
$\omega_m$ are functions of real arguments fixed by
$\sigma_*=\sigma_0+\frac12\tau^2\ell_*$ and $\ell_*$.
\end{proof}

\begin{remark}[the transport reading]\label{rem:reading}
To second order in $\tau$ the deviation is
$\tau^2\bigl(\Delta_1c_1+\frac12\Delta''(0)\bigr)=\frac12\tau^2(\cL''/\cL)'(s_*)+\tau^2\bigl(\alpha'c_1+\frac12\alpha''\bigr)(s_*)$;
the evaluation of Section~\ref{sec:validation} confirms this for $\zeta$.
This deviation is $O(\tau^2)$ uniformly in the height, while the shift from
$w$ to $s_*$ moves the main term by about $\frac14\tau^2K_2\ell_*$. So,
apart from the deviation, the heat flow at height $p$ is evaluation of the
time-zero quantity at $s_*=w+\frac t2\alpha(s_*)$: at height
$p+\frac t2\ell_*$ instead of $p$, and at real part
$x+t\im\alpha(s_*)\approx x+\frac{\pi t}4$. The second shift matches the
drift of the zeros toward the origin at speed about $\pi/4$ that the zero
dynamics $\dot z_j=2\sum_{k\ne j}(z_j-z_k)^{-1}$ give in the mean-field
approximation: the Hilbert transform of the density
$\frac1{4\pi}\log\frac{\lvert x\rvert}{4\pi}$ contributes $-\pi/8$ to the
sum at $x>0$. This is a heuristic and is used in no proof.
\end{remark}

\begin{remark}[the four contributions]\label{rem:costs}
Compared with the time-zero floor, the floor of Theorem~\ref{thm:transport}
gains
$\mathfrak g(t)=\frac12\bigl(\kappa(\sigma_0)-\kappa(\sigma_*)\bigr)$, plus
$\frac14\bigl(\psi(\frac{\sigma_*+\fa}2)-\psi(\frac{\sigma_0+\fa}2)\bigr)$
for the family at $x=0$, and loses $\frac12E$. The loss collects four
contributions: the curvature of $\log\cL$, of order
$\frac12\tau^2(K_2\kappa+\frac12K_3)(\sigma_*)$; the curvature of
$\log\Gamma_F$, of order $\frac12\tau^2(A_1\kappa+\frac12A_2)(\sigma_*)$
with $A_2(\sigma)=-\frac18\psi''(\frac{\sigma+\fa}2)$, which matters for
the family at small height; the Gaussian mass left of $\re s=1$, of order
$\exp(-(\sigma_*-1)^2/2\tau^2)$ times a power of $\mathcal Q_*$; and the
reflected term, of order $e^{-p\ell_*/2}$. As $t\to0$, with
$\ell=2\re\alpha(w)$ the limit of $\ell_*$,
$\mathfrak g/\tau^2\to\frac14K_2(\sigma_0)\ell$ while the first two
contributions divided by $\tau^2$ tend to
$\frac12(k\kappa+\frac12k_2)(\sigma_0)$ with $k=K_2+A_1$ and
$k_2=K_3+A_2$. So at small times the floor does not drop once
$\ell>2(k\kappa+\frac12k_2)/K_2$ at $\sigma_0$. The strip term carries the
Gaussian factor $\exp(-(\sigma_*-\frac12)^2/2\tau^2)$ against a growth of
about $\mathcal Q_*^{1/2}$ at $\sigma=\frac12$. Since
$(\sigma_*-\frac12)^2/2\tau^2=\frac{p^2}{8\tau^2}+\frac{p\ell_*}4+\frac{\tau^2\ell_*^2}8$,
it is exponentially small in $\ell_*$ when $p>2$; for $1<p\le2$ its
smallness comes from the first and last of these terms.
\end{remark}

\begin{remark}[premises]\label{rem:premises}
Theorem~\ref{thm:transport} uses the Hermitian functional equation with a
degree-one Gamma factor; the Euler product, only through
Lemma~\ref{lem:euler} (the coefficients of $\log\cL$ are dominated by those
of $\log\zeta$); bounded coefficients, through Lemma~\ref{lem:partial}; and
about zeros only that $H_t$ has none with $\lvert\im z\rvert>1$.
Linear combinations of $L$-functions of Davenport--Heilbronn type satisfy
the first and third premises and fail the second. The statement
\eqref{eq:Ep} is a smoothing statement at height $p>1$; it constrains
zeros near the axis only through a comparison argument such as
Proposition~\ref{prop:paired}, together with the confinement and
first-contact inputs of Section~\ref{sec:exterior}.
\end{remark}

\section{A bound for \texorpdfstring{$\partial_pS$}{the vertical derivative}}\label{sec:derivative}

The comparison of Section~\ref{sec:strip} also needs an upper bound for
$\partial_pS$. Cauchy's estimate gives one from
Proposition~\ref{prop:E}.

\begin{corollary}[derivative bound]\label{cor:derivative}
Let $r>0$. Suppose that at every point $w'$ of the closed disk
$\lvert w'-w\rvert\le r$ the hypotheses of Lemma~\ref{lem:saddle} and
Proposition~\ref{prop:E} hold, with $E(w')$ the bound \eqref{eq:E}. For the
family this includes $\re w'>1$ and
$\log\frac q\pi+\psi(\frac{\re w'+\fa}2)>0$; for $\zeta$ it includes
$\re w'>1$, $\im w'\ge20$ and $2\re w'-1\le2\im w'$. Let
$a_*\ge\lvert\alpha'(s_*)\rvert$ with $\tau^2a_*<1$, for instance
$a_*=\frac14\psi'(\frac{\sigma_*+\fa}2)$ for the family and $a_*=2/T$ for
$\zeta$ (Lemma~\ref{lem:gamma}). Then
\[
\partial_pS(x,p,t)\ \le\ \frac{K_2(\sigma_*)+a_*}{4(1-\tau^2a_*)}+\frac1{4r}\max_{\lvert w'-w\rvert=r}E(w').
\]
\end{corollary}

\begin{proof}
On the disk, $\mu\mapsto\tau^2\alpha(w'+\mu)$ maps a fixed domain into
itself with a contraction constant uniform in $w'$. For the family the
domain is the half-plane $\re\mu\ge0$. For $\zeta$ it is $D_{x_+}$ of
Lemma~\ref{lem:saddle}, with $x_+$ the largest value of $2\im w'$ on the
disk. The estimates in the proof of that lemma apply to $\mu\in D_{x_+}$,
because there $\lvert s\rvert\le x'+2+\frac18\log x_+\le2x'$ with
$x'=2\im w'$, so the images lie in $D_{x'}\subseteq D_{x_+}$. The iterates
from $\mu=0$
are holomorphic in $w'$ and converge uniformly, so $s_*(w')$ is
holomorphic, and implicit differentiation gives
$s_*'=1/(1-\tau^2\alpha'(s_*))$.

The Gaussian average $\cG$ has no zero with $\re w'>1$
(Lemma~\ref{lem:gauss}). Hence
$\varepsilon(w')=(\cG'/\cG)(w')-(F'/F)(s_*(w'))$ is holomorphic on the
disk, and $\lvert\varepsilon\rvert\le E$ there by
Proposition~\ref{prop:E}. Cauchy's estimate gives
$\lvert\varepsilon'(w)\rvert\le\max_{\lvert w'-w\rvert=r}\lvert\varepsilon(w')\rvert/r$.
Since $S=\frac12\re(\cG'/\cG)(w)$ and $\partial w/\partial p=\frac12$,
\[
\partial_pS=\tfrac14\re\bigl[(F'/F)'(s_*)\,s_*'(w)+\varepsilon'(w)\bigr].
\]
Finally
$\lvert(F'/F)'(s_*)\rvert\le\lvert(\cL'/\cL)'(s_*)\rvert+\lvert\alpha'(s_*)\rvert\le K_2(\sigma_*)+a_*$
by Lemma~\ref{lem:euler} with $j=1$, and
$\lvert s_*'\rvert=\lvert1-\tau^2\alpha'(s_*)\rvert^{-1}\le1/(1-\tau^2a_*)$.
\end{proof}

\begin{remark}\label{rem:branch}
The disk condition matters for $\zeta$ at small heights. There the pole of
$\alpha$ at $s=1$ gives $s_*$ branch points with $\re w$ slightly above
$1$, about $1.055\pm0.625i$ at $t=0.2$ (found by a same-model check, see
Appendix~\ref{app:repro}). Every use below is at $x\ge X$. In the
evaluation, the maximum over the circle is taken at its leftmost point,
height $p-2r$, and at the smallest $\ell_*$; this rests on the sampled
monotonicity of Section~\ref{sec:numerics}, so the evaluated bound is
numerical.
\end{remark}

\section{Numerical evaluation}\label{sec:numerics}

This section evaluates the bound $E$ of Proposition~\ref{prop:E} in floating
point. Three things separate it from a statement for every $t$. First, $E$
is evaluated without outward rounding. Second, the decrease of $E$ in
$\ell_*$, which turns an evaluation at the smallest $\ell_*$ into a bound at
every height, is sampled rather than proved. Third, the times form the
grid
\begin{align*}
\mathcal T=\{&0.001,0.002,0.005,0.01,0.02,0.03,0.05,0.075,0.1,0.125,\\
&0.15,0.175,0.2,0.225,0.25,0.3,0.35,0.4,0.45,0.5\},
\end{align*}
and below $t=0.001$ only the leading order of Remark~\ref{rem:costs} is
controlled. The programmes and records are listed in
Appendix~\ref{app:repro}; R1 is the record of the first programme, with
sections A--H.

\subsection{Method}

$\zeta$ and its derivatives are computed by Euler--Maclaurin summation at
real and complex points, and $\psi,\psi',\psi''$ by recurrence and
asymptotic series; both are checked in R1, section A. The real-variable
majorants are tabulated on logarithmic grids with relative interpolation
error below $4\cdot10^{-5}$, the integrals of Proposition~\ref{prop:E} use
composite Simpson rules, $V=\sigma_*-1-\eta$ with the better of
$\eta\in\{0.03,0.08\}$, and the integrals are cut at $16\tau$ to the right
of $\sigma_*$, at $40\tau$ to the left, and in the reflected region at
$1-\sigma=\sigma_*+\tau^2\ell_*+30\tau+2$. An
independent implementation written by a same-model check (see
Appendix~\ref{app:repro}) differs from the programme by up to
$4\cdot10^{-5}$ relative, in both directions.

\subsection{Thresholds in \texorpdfstring{$\ell_*$}{l*}}

Let $\ell_1(p)$ be the smallest $\ell$ with $\mathfrak g(t)\ge\frac12E$ at
every $t\in\mathcal T$, counting only the $\kappa$ part of the gain
(Remark~\ref{rem:costs}), found by bisection to $0.03$ (R1, section H).
Table~\ref{tab:thresholds} lists it with the small-time threshold
$2(k\kappa+\frac12k_2)/K_2$ of Remark~\ref{rem:costs} at $p=\sqrt5$. For
$\zeta$ the Gamma terms use $T=10^4$; they decrease in $T$, so the row is
valid at every height $x\ge2\cdot10^4+1$, where $\im s_*\ge10^4$.

\begin{table}[ht]
\centering\small
\begin{tabular}{@{}lccccc@{}}
\toprule
Case & $p=\sqrt5$ & $p=2.5$ & $p=3$ & $p=4$ & small-time threshold, $p=\sqrt5$\\
\midrule
family, even & 10.65 & 8.77 & 6.74 & 4.40 & 6.40\\
family, odd & 10.12 & 8.22 & 6.19 & 3.92 & 5.99\\
$\zeta$ ($T=10^4$) & 9.62 & 7.67 & 5.56 & 3.29 & 5.67\\
\bottomrule
\end{tabular}
\caption{The threshold $\ell_1(p)$ (numerical, R1 section H).}\label{tab:thresholds}
\end{table}

Just above $\ell_1$ the binding time is near $t=0.2$: at $\ell=10$ and
$p=\sqrt5$ the even family fails at $t=0.2$ by $0.09$ (R1, section E,
which uses an 11-point grid from $t=0.002$).
From $\ell=11$ on, the margin is smallest at the smallest time of the grid,
where it is $\tau^2$ times the small-time coefficient. With the $\psi$
part of the gain at $x=0$, the small-time thresholds at $p=\sqrt5$ fall to
$5.22$ (even) and $5.38$ (odd). $E$ is nonincreasing in $\ell$, for $\ell$
from 10 to 40, at all 12 sampled settings: the even family at $p=\sqrt5$,
the odd family at $p=3$ and $\zeta$ at $p=\sqrt5$, each at
$t\in\{0.002,0.05,0.2,0.5\}$ (R1, section H). This monotonicity is
checked, not proved, and the other settings used below were not
sampled. With the zeta row, \eqref{eq:Ep}
holds at every grid time from $x\approx4\pi e^{9.62}\approx1.9\cdot10^5$ at
$p=\sqrt5$, and from $x=2\cdot10^4+1$ at $p=3$.

\subsection{The family at every height}\label{sec:family}

At $x=0$ the saddle is real and $\ell_*=\ell_*^{(0)}$, and
Lemma~\ref{lem:saddle} gives $\ell_*(x)\ge\ell_*^{(0)}$ at every $x$. With
$\mathfrak F$ increasing (sampled), $S(x,p,t)\ge\mathfrak F(\ell_*^{(0)}(t))=:\mathfrak F_0(t)$
for every real $x$. Table~\ref{tab:conductor} gives the smallest conductor
for which $\mathfrak F_0(t)\ge C_{q,\fa}(p)$ at every $t\in\mathcal T$:
the passing end of a bisection in $\log_{10}q$ to $0.006$ (R1, section H),
rounded up. These thresholds
lie below those of Table~\ref{tab:thresholds} because the $\psi$ part of
the gain counts at $x=0$.

\begin{table}[ht]
\centering\small
\begin{tabular}{@{}lcc@{}}
\toprule
Parity & $p=\sqrt5$ & $p=3$\\
\midrule
even & $q\ge1.084\cdot10^5$ ($\ell_0=9.50$) & $q\ge1.32\cdot10^3$ ($\ell_0=5.46$)\\
odd & $q\ge4.892\cdot10^4$ ($\ell_0=9.49$) & $q\ge6.62\cdot10^2$ ($\ell_0=5.39$)\\
\bottomrule
\end{tabular}
\caption{Conductors above which $\mathfrak F_0(t)\ge C_{q,\fa}(p)$ on the grid (numerical).}\label{tab:conductor}
\end{table}

At $q=10^6,10^8,10^{10},10^{12}$ the excess
$\min_{t\in\mathcal T}(\mathfrak F_0(t)-C_{q,\fa}(p))$ is positive for both
parities and $p\in\{\sqrt5,3\}$, from $0.0014$ to $0.0077$, attained at
$t=0.001$; at $t=\frac12$ it is $0.37$ to $0.97$ (R1, section F). So on
the grid \eqref{eq:Ep} holds with $C(t)=C_{q,\fa}(p)$ at every real $x$,
and the floor rises with $t$.

\subsection{\texorpdfstring{$\zeta$}{Zeta} beyond the verified height}\label{sec:zetaX}

For $\zeta$ the region of \eqref{eq:Ep} is $x\ge X$. We take Polymath's
height, rounded down to $X=6\cdot10^{10}$ \cite{polymath}, and
$X=5\cdot10^{12}+194858$, the height of the row of Polymath's Table~1 that
Platt and Trudgian use for $\Lambda\le0.2$ \cite{platttrudgian}, as
recorded in \cite{n01}, with
$T=X/2-1$ and $\ell_X=\log\bigl((X/2-1)/2\pi\bigr)-10^{-6}$, which is at most
$\ell_*$ at every $x\ge X$ by Lemma~\ref{lem:saddle}. Table~\ref{tab:zetaX}
gives the excess
$\min_{t\in\mathcal T}\bigl(\mathfrak F(\ell)-(\frac14\ell_X-\frac12\kappa(\sigma_0))\bigr)$,
with $\mathfrak F$ sampled at $\ell_X+\{0,1,3\}$ (R1, section F).

\begin{table}[ht]
\centering\small
\begin{tabular}{@{}lccccc@{}}
\toprule
$X$ & $p=\sqrt5$ & $p=3$ & $p=4$ & $p=5$ & $p=6$\\
\midrule
$6\cdot10^{10}$ ($\ell_X=22.29$) & 0.0051 & 0.0021 & 0.00088 & 0.00044 & 0.00025\\
$5\cdot10^{12}+194858$ ($\ell_X=26.71$) & 0.0064 & 0.0026 & 0.0011 & 0.00054 & 0.00030\\
\bottomrule
\end{tabular}
\caption{Excess of the floor over $\frac14\ell_X-\frac12\kappa(\sigma_0)$ for $\zeta$ at $x\ge X$ (numerical).}\label{tab:zetaX}
\end{table}

All minima are attained at $t=0.001$. They are thin at large $p$ only
because $\tau^2$ is small there, so the floor at those times is close to
$\frac14\ell_X-\frac12\kappa(\sigma_0)$. On the grid, then, \eqref{eq:Ep}
holds at every $x\ge X$ with the constant floor
$\frac14\ell_X-\frac12\kappa(\sigma_0)$; the $x$-dependent floor
$\frac14\log\frac x{4\pi}-\frac12\kappa(\sigma_0)$ is larger. The sampled
heights include the band $3\le\im z\le6$ on which \cite{hb} works.

\subsection{Validation}\label{sec:validation}

\emph{The Gaussian representation} (R1, section B). For $\zeta$ at
$(x,p,t)=(3,\sqrt5,0)$, $(10,\sqrt5,0.1)$, $(20,3,0.3)$ and $(25,\sqrt5,0.5)$,
$S$ from Polymath's kernel integral and $S$ from the Gaussian average of
$\xi$ agree to $9\cdot10^{-16}$. The two routes share no formula beyond the
definition of $\xi$.

\emph{The transport law} (R1, sections C and D). For $\zeta$ at $x=10^3$
and $10^4$ ($\ell_*=4.38$ and $6.68$), $p\in\{\sqrt5,3\}$ and
$t\in\{0.02,0.1,0.3,0.5\}$, we compare $S$ from the Gaussian average of
$\xi$ with the main term $\frac12\re(F'/F)(s_*)$ and the second-order term
of Remark~\ref{rem:reading}. For characters (real characters modulo 10007,
odd; 10009, even; 100003 and 1000003, odd; 1000033, even; and a cubic
character modulo 10009) at $x\in\{0,2,10\}$, $p\in\{\sqrt5,3\}$ and
$t\in\{0.05,0.2,0.5\}$, $S$ is computed from the kernel integral with
Gauss-sum root numbers, and the main term as the kernel's $S$ at time zero
at $(2\im s_*,2\re s_*-1)$. The findings are as follows.
\begin{itemize}
\item The deviation $\lvert\varepsilon\rvert$ (for $\zeta$, of
$\cG'/\cG$; for characters, of its real part) is at most $0.044$ for
$\zeta$ and $0.117$ for the characters. The largest is at $x=0$ for the
real character modulo 10009 at $t=\frac12$.
\item For $\zeta$ the second-order term accounts for most of the
deviation: at $x=10^3$, $p=\sqrt5$ and $t=0.02$ it is $0.0047$ against a
measured $0.0048$.
\item $E\ge\lvert\varepsilon\rvert$ at all 124 points. The ratio
$E/\lvert\varepsilon\rvert$ runs from $1.6$ (even character modulo
1000033, $x=0$, $p=3$, $t=\frac12$) to about 1300, with median 22. It is
largest at $p=\sqrt5$ and moderate $\ell_*$.
\item $S(x,p,t)$ exceeds the time-zero floor at the same $x$ at every
point, including points below the thresholds of
Table~\ref{tab:thresholds}.
\end{itemize}
These points test the identity and the size of $E$ at moderate heights.
They do not test the regime of Section~\ref{sec:zetaX}, where the theorem,
not a computation, is the evidence.

\subsection{Comparison with Hakimoglu-Brown's floor}\label{sec:compare}

Hakimoglu-Brown's Euler-product floors for $S$ at height $p$
\cite[Propositions~8 and~11]{hb} have the same main term as
Theorem~\ref{thm:transport},
$\frac14L-\frac12\kappa(\frac{1+p}2+\frac t4L)$ with $L=\log\frac x{4\pi}$.
Two identities explain this. His heat weight is the Gaussian average,
\[
e^{t\log^2n/4}=\E\,n^{-\tau N}\qquad(\tau^2=t/2),
\]
and his point $s_B=w_B+\frac t2\alpha_B(w_B)$, with $w_B=\frac12(1-iz)=\bar w$
and $\alpha_B(s)=\frac1{2s}+\frac1{s-1}+\frac12\log\frac s{2\pi}$, is the
complex conjugate of the first iterate of the saddle equation
$s=w+\frac t2\alpha(s)$, with $\alpha$ in Stirling form;
$\alpha-\alpha_B=-\frac1{6s^2}+O(s^{-4})$. He controls the remainder
through Polymath's Riemann--Siegel coefficient bounds
\cite[Propositions~6.1--6.2]{polymath} on the band $3\le\im z\le6$ and a
truncated, heat-weighted Dirichlet polynomial. Here it is controlled by
the exact Gaussian representation and Euler-product majorants, at any
$p>1$.

Table~\ref{tab:compare} compares the two at his height $p=61/20$ and his
$X=60000083951.5$, on intervals $[t^-,t^-+\frac1{200}]$ (R2, section B).
His bounds are his expressions (31)--(32) with his own upper-bound inputs,
which reproduce his certified Table~3 (R2, section A). Ours take the gain
at $t^-$ and $E$ at $t^+$, with $G$ from Corollary~\ref{cor:derivative}
and the best radius $r\in\{0.05,\dots,0.35\}$.

\begin{table}[ht]
\centering\small
\begin{tabular}{@{}ccccccc@{}}
\toprule
$t^-$ & main term & his $C^*$ & our $C$ & his $G^*$ & our $G$ & main term of $G$\\
\midrule
0 & 5.2974 & 5.2903 & 5.2958 & 0.2340 & 0.2174 & 0.2094\\
0.05 & 5.3867 & 5.3603 & 5.3777 & 0.1738 & 0.1601 & 0.1217\\
0.10 & 5.4408 & 5.4193 & 5.4316 & 0.1121 & 0.1111 & 0.0771\\
0.15 & 5.4761 & 5.4607 & 5.4683 & 0.0735 & 0.0769 & 0.0517\\
0.195 & 5.4982 & 5.4869 & 5.4917 & 0.0520 & 0.0562 & 0.0374\\
\bottomrule
\end{tabular}
\caption{Floors for $S$ and bounds for $\partial_pS$ at $p=61/20$, $x\ge X$
(floating; his values reproduce his certified ones).}\label{tab:compare}
\end{table}

The loss against the main term is $0.007$--$0.026$ for his floor and
$0.0016$--$0.0095$ for ours, so our remainder is $1.7$ to $4.4$ times
smaller. Our derivative bound is smaller for $t\le0.1$ and up to $8\%$
larger after. At $p=301/100$, against his Table~2 (Proposition~8, steps of
$\frac1{100}$), our floor is larger by $0.012$ to $0.025$. These
comparisons set a floating evaluation against a certified one: the
advantage counts only once our evaluation is certified.

\section{A check from the classical strip}\label{sec:strip}

This section restarts Hakimoglu-Brown's derivative-kernel comparison
\cite[Theorem~12]{hb} from the classical strip $M(0)\le1$ instead of his
hypothesis $M(0)\le\frac9{16}$. Its analytic inputs are those of \cite{hb},
an unrefereed working draft. We use them as stated and do not re-derive
them. Throughout, $X=60000083951.5$, $b=\frac1{25}$ and $p=\frac{61}{20}$,
so $B=p^2=9.3025$; the squared heights that \cite{hb} calls $q_i$ are
written $Q_i$.

\subsection{Imported statements}\label{sec:imported}

The four statements below are quoted from \cite{hb}, with its notation
adapted to ours. They are not proved here.

\begin{lemma}[{\cite[Lemma~2]{hb}}; imported]\label{lem:hb2}
For $0\le t\le\frac15$, a zero of $H_t$ with $\lvert\re z\rvert\le X$
satisfies $\lvert\im z\rvert<\frac15$.
\end{lemma}

Its proof uses Polymath's boundary computation
\cite[Section~8.4, claim~(a)]{polymath}, that $H_t(x+iy)\ne0$ for
$X\le x\le X+1$, $\frac15\le y\le1$ and $0\le t\le\frac15$; the verification
of the Riemann hypothesis below height $\frac12(X+1)$; backward Hermite
splitting \cite[Proposition~3.1(ii)]{polymath}; and the classical strip
theorem.

\begin{lemma}[{\cite[Lemma~10]{hb}}; imported]\label{lem:hb10}
Let $p>0$, $B=p^2$, $0<A_2<B$ and $\lambda,\mu\ge0$. For $D\ge0$ and
$0\le r\le A_2$ put $w=D+B-r$, $\Delta_{\rm K}=w^2+4Dr$ and
\begin{equation}\label{eq:Pi}
\Pi(D,r)=2\Delta_{\rm K}^2-(D+A_2+3r)\Bigl(2\lambda w\Delta_{\rm K}-\mu(2w+4B)\Delta_{\rm K}+8\mu Bw^2\Bigr).
\end{equation}
If $\Pi\ge0$ throughout that domain, then for $0<h^2\le A_2$, $0\le v\le h$
and $(d,v)\ne(0,h)$,
\[
K_h(d,v)\ \ge\ h\Bigl(\frac\lambda pK_p(d,v)-\mu\,\partial_pK_p(d,v)\Bigr).
\]
\end{lemma}

For $\lambda=1$ and $\mu=0$ this is the kernel inequality in the proof of
Proposition~\ref{prop:paired}; the derivative term pays for $\lambda>1$.
We rechecked the proof of Lemma~\ref{lem:hb10}; it holds for every
$A_2<B$.

\begin{proposition}[{\cite[Proposition~11]{hb}}; imported]\label{prop:hb11}
Let $p=\frac{61}{20}$ and $0\le t^-\le t^+\le\frac15$. Put
$\sigma_p=\frac{1+p}2=\frac{81}{40}$, $u=\sigma_p+5.57\,t^-$,
$u_8=\sigma_p+2.78\,t^-$, $\bar e=\frac{p-1}2-2.78\,t^-$, $N_1=65536$,
$N_*=69098$, $L_*=2\log N_*+3$, and for $j=0,1,2$
\begin{gather*}
\delta_j=\sum_{n=2}^{N_1}\bigl(e^{t^+\log^2n/4}-1\bigr)(\log n)^jn^{-u}+\int_{N_1}^\infty(\log y)^jy^{-u_8}\dd y,\\
\varrho_0=1.01\,N_*^{-p}\Bigl(\frac{N_*^{\bar e+1}}{\bar e+1}+N_*^{\bar e}\Bigr),\qquad\varrho_1=L_*\varrho_0,\qquad\varrho_2=(L_*^2+10^{-9})\varrho_0,\\
\eta_0=10^{-6},\quad\eta_1=4\cdot10^{-5},\quad\eta_2=1.6\cdot10^{-3},\quad\nu_1=\tfrac12+10^{-10},\quad\nu_2=\tfrac14+10^{-10},\\
e_0=\delta_0+\varrho_0+\eta_0,\qquad e_1=\nu_1\delta_1+\varrho_1+\eta_1,\qquad e_2=\nu_2\delta_2+\varrho_2+\eta_2+10^{-9},\\
\varepsilon_0=\zeta(u)e_0,\quad\varepsilon_1=\zeta(u)\bigl(e_1+\nu_1\kappa(u)e_0\bigr),\quad
\varepsilon_2=\zeta(u)\bigl\{e_2+2\nu_1\kappa(u)e_1+\nu_2\bigl(K_2(u)+\kappa(u)^2\bigr)e_0\bigr\}.
\end{gather*}
If $\varepsilon_0<\frac12$, then for $x\ge X$ and $t^-\le t\le t^+$
\begin{align*}
S(x,p,t)&\ \ge\ C^*:=5.5716-10^{-6}-\nu_1\kappa(u)-\frac{\varepsilon_1}{1-\varepsilon_0},\\
\partial_pS(x,p,t)&\ \le\ G^*:=\nu_2K_2(u)+\frac{\varepsilon_2}{1-\varepsilon_0}+\Bigl(\frac{\varepsilon_1}{1-\varepsilon_0}\Bigr)^2+2\cdot10^{-10}.
\end{align*}
\end{proposition}

These are \cite[(28)--(32)]{hb}, where $\zeta(u)$, $\kappa(u)$, $K_2(u)$,
$\nu_j$, $\varrho_j$ and $\bar e$ are written $Z$, $\psi$, $\psi_2$, $c_j$,
$\rho_j$ and $e$. The proof uses \cite[Lemma~3]{hb}, an approximation of
Riemann--Siegel type on $3\le\im z\le6$ derived from Polymath's
decomposition \cite[(37)--(41)]{polymath} and coefficient bounds
\cite[Propositions~6.1--6.2]{polymath}; its Lemma~6; the normalizer bounds
of its Section~3; its fixed-cutoff argument around~(16); its
estimate~(27); and the elementary inequalities of its Section~8.2, among
them $\frac14\log\frac x{4\pi}>5.5716$ and
$\lfloor\sqrt{x/4\pi}\rfloor\ge N_*$ for $x\ge X$.

\begin{theorem}[{\cite[Theorem~12]{hb}}; imported]\label{thm:hb12}
Assume $M(0)\le Q_0<B=p^2$. Choose $0=t_0<\dots<t_n\le\frac15$,
$d_i=t_{i+1}-t_i$, and $Q_0>Q_1>\dots>Q_n\ge b$. For $0\le i<n$ suppose
that $C_i,G_i,\lambda_i,\mu_i\ge0$ satisfy
\[
S(x,p,t)\ge C_i,\qquad\partial_pS(x,p,t)\le G_i\qquad(x\ge X,\ t_i\le t\le t_{i+1}),
\]
that \eqref{eq:Pi} with $(A_2,\lambda,\mu)=(Q_i,\lambda_i,\mu_i)$ is
nonnegative for $D\ge0$ and $0\le r\le Q_i$, and that
\[
k_i=4\Bigl\{\frac{\lambda_iC_i}p-\mu_iG_i-\frac{2\lambda_i}{B-Q_i}-\frac{2\mu_i(B+Q_i)}{(B-Q_i)^2}\Bigr\}>0 .
\]
If $e^{k_id_i}(2+k_iQ_{i+1})\ge2+k_iQ_i$ for every $i$, then
$M(t_n)\le Q_n$ and $\Lambda\le t_n+\frac12Q_n$.
\end{theorem}

Its proof uses Lemmas~\ref{lem:hb2} and~\ref{lem:hb10}, the zero dynamics
of Proposition~\ref{prop:paired}, the classical strip theorem, backward
Hermite splitting, and \cite[Theorem~1.5(i)]{polymath}, which places
every nonreal zero at $\lvert x\rvert<\exp(C/t)$ with an absolute constant
$C$, for compact attainment of a first contact.

\subsection{Where the initial strip enters}\label{sec:q0}

Hakimoglu-Brown states that his hypothesis ``fixes the initial value
$q_0=9/16$ of the strip comparison in Theorem~12; the zero dynamics, the
boundary input, the approximation on $3\le y\le6$, the kernel inequality,
and all certificates are independent of it'' \cite[Section~1.1]{hb}, and
he calls Theorem~12 ``a general conversion theorem''. We checked this
against the proofs. Lemma~\ref{lem:hb2} uses the classical bound
$\lvert\im z\rvert\le1$, not $Q_0$. The barrier argument of
Theorem~\ref{thm:hb12} needs the barrier to start above $M(0)$, together
with the inputs listed after it. Lemma~\ref{lem:hb10} needs only
$A_2<B$. With $Q_0=1$ we have $M(0)\le1$, $Q_0<B$, and a final time below
$\frac15$. His Section~1.1 lists other half-planes only for $\sigma<1$, and
his Remark~5 and Section~9 warn against extrapolating his closed one-step
formulas to $\sigma=1$. Those warnings concern only the closed formulas.
Indeed, from $Q_0=1$ his closed formulas do not reach $\frac1{25}$ by
$t=\frac15$ \cite[Remark~5]{hb}, and neither does the comparison of
Proposition~\ref{prop:paired} with the floors of
Proposition~\ref{prop:hb11} (R2, section C).

\begin{proposition}[conditional on \cite{hb}]\label{prop:table}
Assume Lemmas~\ref{lem:hb2} and~\ref{lem:hb10},
Proposition~\ref{prop:hb11} and Theorem~\ref{thm:hb12}, together with the
inputs their proofs use: from \cite{hb}, Lemmas~3 and~6, the normalizer
bounds of Section~3, the fixed-cutoff argument around~(16), estimate~(27)
and the elementary inequalities of Section~8.2; from \cite{polymath}, the
boundary computation of Section~8.4, equations (37)--(41),
Propositions~3.1(ii), 6.1 and~6.2 and Theorem~1.5(i); the verification of
the Riemann hypothesis to height $\frac12(X+1)$ \cite{platt2017}; and the
classical strip theorem. Then
\[
\Lambda\ \le\ \frac{18891}{100000}=0.18891 .
\]
\end{proposition}

\begin{proof}
Apply Theorem~\ref{thm:hb12} with $Q_0=1$, $n=34$, $t_i=i/200$ for
$i\le33$, $d_i=\frac1{200}$ for $i<33$ and $d_{33}=\frac{391}{100000}$, so
that $t_{34}=\frac{16891}{100000}$, with $Q_{34}=b$ and the rows
$(Q_i,C_i,G_i,\lambda_i,\mu_i)$ of Table~\ref{tab:table} in
Appendix~\ref{app:table}. The classical strip gives $M(0)\le1=Q_0<B$. The
certificate R3 (Appendix~\ref{app:repro}) verifies the remaining
hypotheses.
\begin{itemize}
\item In exact rational arithmetic: the chain ($t_{i+1}=t_i+d_i$,
$Q_{i+1}<Q_i$, $Q_{i+1}\ge b$, $Q_{34}=b$, $t_{34}\le\frac15$; the last
step $Q_{34}<Q_{33}$ is read off the table); the
nonnegativity of \eqref{eq:Pi} on every row, through the positivity of all
$34\cdot8\cdot25=6800$ Bernstein coefficients of the compactification
\cite[(34)--(35)]{hb} on eight subintervals, the smallest being
$1.4155\cdot10^{-3}$; $7.780<k_i<10.592$; every gate, with the
exponential replaced by its Taylor sum $\mathrm T_{20}(y)=\sum_{j\le20}y^j/j!\le e^y$,
smallest margin $1.079\cdot10^{-5}$; and the two inequalities of
\cite[Section~8.2]{hb} displayed after Proposition~\ref{prop:hb11}. The
other elementary inequalities of that section are taken from \cite{hb}.
\item In directed-rounding decimal arithmetic at 32 digits, with $\exp$ and
$\log$ moved one unit outward and the other operations rounded toward the
bound: on every row, $C_i\le C^*$ and $G_i\ge G^*$ for
Proposition~\ref{prop:hb11} on $[t_i,t_i+\frac1{200}]\supseteq[t_i,t_{i+1}]$,
with margins at least $2.05\cdot10^{-6}$ and $2.02\cdot10^{-6}$, and
$\varepsilon_0\le0.0191<\frac12$. Here $\zeta(u)$, $\kappa(u)$, $K_2(u)$ and
the $\delta_j$ are replaced by the finite upper bounds of
\cite[Section~8.2]{hb}; $C^*$ decreases and $G^*$ and $\varepsilon_0$
increase in these inputs.
\end{itemize}
Theorem~\ref{thm:hb12} then gives $M(t_{34})\le\frac1{25}$ and
$\Lambda\le t_{34}+\frac1{50}=\frac{18891}{100000}$.
\end{proof}

\subsection{Control, construction and sensitivity}

\emph{Control.} The same programme, applied to \cite[Table~3]{hb}
($Q_0=\frac9{16}$), returns its printed certificates: the least Bernstein
coefficient
$623287907570432169922330885087/479268282493202500000000000000000$ in row
21, a gate margin above $2.78\cdot10^{-8}$, scalar margins
$2.05\cdot10^{-6}$ and $1.10\cdot10^{-6}$, $8.63<k_i<10.32$ and the final
time $\frac{11963}{100000}$.

\emph{Construction} (R2, section C). The $C_i$ and $G_i$ are
Proposition~\ref{prop:hb11} on $[t_i,t_i+\frac1{200}]$, rounded down and up
to $10^{-4}$ with a margin of at least $2\cdot10^{-6}$; they do not depend
on $Q$. Rows 0--23 coincide with
\cite[Table~3]{hb}, except that $G_6$ is $10^{-4}$ larger here; his row 23
covers $[0.115,0.11963]$ and ours $[0.115,0.12]$. The $\lambda_i,\mu_i$
maximize $k_i$ over a sampled admissible region, are reduced by $0.005$,
and were then rounded and lowered until the exact Bernstein test passed.
$Q_{i+1}$ is the barrier value $(Q_i+2/k_i)e^{-k_id_i}-2/k_i$ rounded up
to $10^{-6}$, plus $10^{-6}$.

\emph{Sensitivity} (floating, R2 section C). A step of $\frac1{400}$ gives
$0.18850$, a step of $\frac1{100}$ gives $0.18961$, and a polynomial margin
of $0.02$ gives $0.18982$.

\emph{Thin places.} The scalar margins of about $2\cdot10^{-6}$ are
Hakimoglu-Brown's own; since rows 0--23 share his $C_i$ and, except
$G_6$, his $G_i$, a correction to his constants would affect both tables
alike. In row 0 the leading coefficient of \eqref{eq:Pi} in $D$,
$2-2\lambda_0+2\mu_0=0.014$, is exact and positive. The smallest gate
margin is $1.08\cdot10^{-5}$.

\subsection{Literature}\label{sec:literature}

A same-model literature search, made after the table was built, read what
it could fetch through a summarizing tool. Its quotations were not checked
against the full texts, so Table~\ref{tab:literature} reports the claims as
their authors state them.

\begin{table}[ht]
\centering\small
\begin{tabularx}{\textwidth}{@{}>{\raggedright\arraybackslash}p{0.2\textwidth}>{\raggedright\arraybackslash}p{0.3\textwidth}X@{}}
\toprule
Source & Claim and status & Relation to Proposition~\ref{prop:table}\\
\midrule
Gordon \cite{gordon} & $\Lambda\le0.157656619$, no half-plane hypothesis; ``has not undergone external peer review''; uses Platt and Trudgian's verification to $3\cdot10^{12}$ & dynamics of the highest zero with a reference barrier started at $(1000001/1000000)^2$, ``the global strip gives strict entry at zero''; own analytic inputs and a far larger $X$\\
Raymond \cite{raymond} & $\Lambda<0.149396$; ``has not undergone external peer review'' & ``starts from the classical width-one zero strip at time zero''; 608 time cells\\
Planat \cite{planat} & $\Lambda\le0.148611789321711$ under stated premises & ``a directed comparison beginning with the classical zero strip'', using Raymond's source data; its appendix on \cite{hb} was not read\\
Gomila \cite{gomila} & $\Lambda\le0.1787854$, ``not yet been peer reviewed'' & Polymath's criterion at $X\approx6\cdot10^{12}$ with a barrier computation; no strip comparison\\
Raymond, conditional \cite{raymondcond} & $\Lambda<0.116398$ assuming no zero with $\re s>\frac78$ & below \cite{hb}'s $0.13963$ under the same hypothesis\\
\bottomrule
\end{tabularx}
\caption{Unrefereed 2026 claims found by the literature search.}\label{tab:literature}
\end{table}

No later version of \cite{hb} was found, and nothing stating $0.18891$ or
restarting its Theorem~12 from $Q_0=1$ turned up. Three imported inputs
were checked. Platt \cite{platt2017} isolates the zeros of $\zeta$ with
imaginary part up to $30{,}610{,}046{,}000$, which exceeds
$\frac12(X+1)=30{,}000{,}041{,}976.25$ by about $6\cdot10^8$. Polymath's
Section~1 records the barrier region $X\le x\le X+1$, $y_0\le y\le1$,
$0\le t\le t_0$ with $t_0=0.2$, $X=6\times10^{10}+83952-0.5$ and $y_0=0.2$
\cite{polymath}; its Section~8.4 could not be fetched in full. Polymath's
Theorem~1.5 holds for $0<t\le\frac12$ with an absolute, non-explicit
constant, as the barrier argument of Theorem~\ref{thm:hb12} needs.

Starting a stepped comparison from the classical strip is therefore not new,
and Proposition~\ref{prop:table} is not a record. Its distinctive feature
is that its numerical inputs are published, refereed computations,
Polymath's barrier and Platt's verification, with the
analytic lemmas of \cite{hb} on top, checked by a 34-row certificate. The
stronger claims rest on their authors' own unrefereed computations. For
this paper the zeta value is a cross-check: with published numerical inputs
only, a stepped comparison started from the classical strip already goes
below the published $0.2$, consistent with those claims.

\section{Planning: what the floor would give}\label{sec:planning}

The values in this section are planning arithmetic. They combine the
floating evaluation of Section~\ref{sec:numerics} with comparison
arguments, some of which are not yet written, and they bound nothing.

\subsection{\texorpdfstring{$\zeta$}{Zeta} with the transport floor}

Replace Proposition~\ref{prop:hb11} by Theorem~\ref{thm:transport} and
Corollary~\ref{cor:derivative}, evaluated as in Section~\ref{sec:compare},
and run Theorem~\ref{thm:hb12}, or the comparison of
Proposition~\ref{prop:paired} (\cite[Theorem~9]{hb}), at
$X=60000083951.5$ inside the window $t\le\frac15$ of Lemma~\ref{lem:hb2},
with steps of $\frac1{200}$ (R2, section D). Table~\ref{tab:zetaplan}
gives the resulting planning values for $\Lambda$; a dash means that
$\frac1{25}$ is not reached by $t=\frac15$, or that $p^2<5Q_0$.

\begin{table}[ht]
\centering\small
\begin{tabular}{@{}ccccc@{}}
\toprule
$p$ & Thm.~\ref{thm:hb12}, $Q_0=1$ & Thm.~\ref{thm:hb12}, $Q_0=\frac9{16}$ & Prop.~\ref{prop:paired}, $Q_0=1$ & Prop.~\ref{prop:paired}, $Q_0=\frac9{16}$\\
\midrule
2.0 & 0.1953 & 0.1441 & -- & 0.1491\\
2.25 & 0.1885 & 0.1382 & 0.2085 & 0.1524\\
2.5 & 0.1869 & \textbf{0.1370} & 0.2148 & 0.1561\\
2.75 & \textbf{0.1867} & 0.1374 & -- & 0.1599\\
3.05 & 0.1882 & 0.1391 & -- & 0.1643\\
3.5 & 0.1934 & 0.1424 & -- & 0.1707\\
\bottomrule
\end{tabular}
\caption{Planning values for $\zeta$ with the floor of Theorem~\ref{thm:transport} (R2, section D).}\label{tab:zetaplan}
\end{table}

So the transport floor would move $0.18891$ and Hakimoglu-Brown's
conditional $0.13963$ by about $0.002$--$0.003$. Heights below $2.5$ do not
help, for two reasons: the polynomial of Lemma~\ref{lem:hb10} leaves less
room for $\lambda>1$ when $B=p^2$ is small, and $K_2$ grows as $\sigma$
approaches $1$, which raises $G$. Once certified,
Theorem~\ref{thm:transport} and Corollary~\ref{cor:derivative} would
supply the inputs of such a table without the Riemann--Siegel
approximation of \cite[Lemma~3]{hb}. The comparison itself would still
rest on Lemmas~\ref{lem:hb2} and~\ref{lem:hb10} and
Theorem~\ref{thm:hb12}.

\subsection{The Dirichlet family}

For Dirichlet $L$-functions the transport floor is the main new input. Take
every primitive character of conductor $q$, start from $Q_0=1$, and let
$b\to0$: by Section~\ref{sec:family}, \eqref{eq:Ep} then holds at every
real $x$ and no finite region is needed. Table~\ref{tab:familyplan} gives
the time at which the barrier reaches $0$, for the best $p$ on a small
grid. The first column uses the constant floor $C_{q,\fa}(p)$ with
Proposition~\ref{prop:paired} \cite[Section~4.6]{p01}; the second uses
$\mathfrak F_0(t)$ with the same comparison (R1, section G); the third uses
$\mathfrak F_0(t)$ and Corollary~\ref{cor:derivative} with the kernel of
Lemma~\ref{lem:hb10} (R2, section E). R2's own run with the kernel of
Proposition~\ref{prop:paired} agrees with the second column to about
$0.001$.

\begin{table}[ht]
\centering\small
\begin{tabular}{@{}cccccc@{}}
\toprule
$q$ & parity & constant floor & transport floor & transport floor, Lemma~\ref{lem:hb10} kernel & change\\
\midrule
$10^6$ & even & 0.3161 (2.85) & 0.3076 (2.70) & 0.2920 (4.0) & $-7.6\%$\\
$10^6$ & odd & 0.3073 (2.69) & 0.2990 (2.54) & 0.2827 (3.5) & $-8.0\%$\\
$10^8$ & even & 0.2618 (2.37) & 0.2519 (2.37) & 0.2344 (3.0) & $-10.5\%$\\
$10^8$ & odd & 0.2541 (2.27) & 0.2447 (2.27) & 0.2268 (3.0) & $-10.7\%$\\
$10^{12}$ & even & 0.1936 (2.24) & 0.1868 (2.24) & 0.1651 (2.5) & $-14.7\%$\\
$10^{12}$ & odd & 0.1897 (2.24) & 0.1839 (2.24) & 0.1611 (2.5) & $-15.1\%$\\
\bottomrule
\end{tabular}
\caption{Planning values for $\Lambda_\chi$, uniform over primitive characters of conductor $q$, with the best $p$ in parentheses. The last column compares the third column with the first, using the printed values.}\label{tab:familyplan}
\end{table}

The best $p$ for the even character at $q=10^6$ is at the edge of the grid.
These values remain conditional on the uniform confinement of the nonreal
zeros at positive times (\cite{p01}, item H12), on a written comparison
theorem for the family, now with Lemma~\ref{lem:hb10}'s kernel and the
second exterior input $\partial_pS\le G$, and on the certification of
Section~\ref{sec:numerics}.

For a single character, Platt's verification heights \cite{platt2016} give
$\ell\approx16.6$ for every $q\le400{,}000$ \cite[Section~4.1]{p01}, above
every $\ell_1$ of Table~\ref{tab:thresholds}. So on the grid
\eqref{eq:Ep} holds with $C(t)=\frac14\ell-\frac12\kappa(\sigma_0)$ beyond
the verified height of each such character. A bound for $\Lambda_\chi$
would still need a barrier for that character and a finite-region argument.

\subsection{What remains}

\begin{enumerate}
\item \emph{Review.} Sections~\ref{sec:setting}--\ref{sec:derivative}
should be reviewed first, since everything else rests on
Theorem~\ref{thm:transport} and Corollary~\ref{cor:derivative}.
\item \emph{Certification.} Re-evaluate $E$ and the bound of
Corollary~\ref{cor:derivative} with outward rounding, with monotone
majorants on intervals of $\sigma$ and interval quadrature; cover every
$t\in(0,\frac12]$ by monotonicity or Lipschitz bounds between grid times;
prove the decrease of $E$ in $\ell_*$ term by term; and bound the
$O(\tau^3)$ term below $t=0.001$. If outward rounding moves
$\ell_1(\sqrt5)$ above 12, the conductor thresholds of
Table~\ref{tab:conductor} must be recomputed.
\item \emph{The family.} Confinement (H12) and the comparison theorem
\cite{p01}.
\item \emph{$\zeta$.} An independent re-derivation of \cite[Lemma~3 and
Proposition~11]{hb}, on which Proposition~\ref{prop:table} rests.
\item \emph{Literature.} Full readings of
\cite{gordon,raymond,planat,gomila}.
\end{enumerate}

\clearpage
\appendix

\section{The table of Proposition~\ref{prop:table}}\label{app:table}

Here $p=\frac{61}{20}$, $t_i=i/200$, $d_i=\frac1{200}$ for $i<33$,
$d_{33}=\frac{391}{100000}$ and $Q_{34}=\frac1{25}$. Each next height is the
following row's $Q_i$. The columns are $10^6Q_i$, $10^4C_i$, $10^4G_i$,
$10^3\lambda_i$ and $10^3\mu_i$; all entries are exact.

\begin{table}[h!]
\centering\small
\begin{tabular}{@{}rrrrrr@{\qquad}rrrrrr@{}}
\toprule
$i$ & $10^6Q_i$ & $10^4C_i$ & $10^4G_i$ & $10^3\lambda_i$ & $10^3\mu_i$ & $i$ & $10^6Q_i$ & $10^4C_i$ & $10^4G_i$ & $10^3\lambda_i$ & $10^3\mu_i$\\
\midrule
0 & 1000000 & 52902 & 2340 & 1474 & 481 & 17 & 356111 & 54036 & 1280 & 1804 & 1068\\
1 & 952036 & 52964 & 2322 & 1524 & 540 & 18 & 329670 & 54090 & 1225 & 1815 & 1096\\
2 & 904955 & 53031 & 2284 & 1574 & 609 & 19 & 304348 & 54143 & 1172 & 1824 & 1118\\
3 & 858867 & 53100 & 2234 & 1621 & 691 & 20 & 280116 & 54193 & 1122 & 1834 & 1139\\
4 & 813966 & 53172 & 2174 & 1644 & 738 & 21 & 256931 & 54242 & 1074 & 1844 & 1167\\
5 & 770525 & 53245 & 2107 & 1658 & 764 & 22 & 234774 & 54288 & 1028 & 1853 & 1188\\
6 & 728552 & 53318 & 2036 & 1671 & 788 & 23 & 213607 & 54334 & 985 & 1861 & 1208\\
7 & 688021 & 53390 & 1962 & 1684 & 815 & 24 & 193403 & 54377 & 944 & 1869 & 1227\\
8 & 648917 & 53462 & 1887 & 1698 & 841 & 25 & 174128 & 54419 & 905 & 1878 & 1255\\
9 & 611181 & 53533 & 1813 & 1710 & 865 & 26 & 155754 & 54459 & 867 & 1885 & 1273\\
10 & 574812 & 53602 & 1739 & 1723 & 892 & 27 & 138247 & 54498 & 832 & 1892 & 1291\\
11 & 539777 & 53670 & 1667 & 1735 & 916 & 28 & 121575 & 54536 & 798 & 1899 & 1308\\
12 & 506049 & 53736 & 1596 & 1748 & 944 & 29 & 105705 & 54572 & 766 & 1905 & 1325\\
13 & 473600 & 53800 & 1528 & 1759 & 967 & 30 & 90610 & 54606 & 736 & 1912 & 1341\\
14 & 442411 & 53862 & 1462 & 1771 & 995 & 31 & 76255 & 54640 & 707 & 1919 & 1367\\
15 & 412455 & 53922 & 1399 & 1782 & 1018 & 32 & 62616 & 54672 & 679 & 1925 & 1383\\
16 & 383698 & 53980 & 1338 & 1794 & 1046 & 33 & 49660 & 54703 & 653 & 1930 & 1397\\
\bottomrule
\end{tabular}
\caption{The 34 rows from $Q_0=1$ (certificate R3).}\label{tab:table}
\end{table}

\section{Reproducibility and checks}\label{app:repro}

\subsection{Programmes and records}

The programmes and records are in
\nolinkurl{papers/investigations-3/03_effective_outer_remainders/numerics/round-1/}
of the repository \cite{records}. They are standard-library Python and
each prints one deterministic JSON record. Run each from that folder as
\begin{center}
\texttt{python3 -I <programme>.py > <RECORD>.json}
\end{center}
\begin{itemize}
\item[P1] \nolinkurl{height_p_transport_claude_20261017.py}, record R1
\nolinkurl{HEIGHT_P_TRANSPORT_CLAUDE_20261017.json} (106\,kB). Sections A--D
are floating diagnostics, E, F and H a floating evaluation of
Proposition~\ref{prop:E}, and G planning. A few minutes.
\item[P2] \nolinkurl{hakimoglu_brown_comparison_claude_20261017.py}, record
R2 \nolinkurl{HAKIMOGLU_BROWN_COMPARISON_CLAUDE_20261017.json} (19\,kB). It
imports P1. Section A reproduces printed values of \cite{hb}: its value
(18), the gates of its Table~2, the margins of its Proposition~8, the
margins of its Table~3 in floating point, and its Tables~6--7 at rows 0, 12
and 23. Section B is the floating
comparison of Section~\ref{sec:compare}; C builds the table of
Appendix~\ref{app:table}, with exact Bernstein and gate checks and floating
scalars; D and E are planning. A few minutes.
\item[P3] \nolinkurl{unconditional_strip_table_claude_20261017.py}, record
R3 \nolinkurl{UNCONDITIONAL_STRIP_TABLE_CLAUDE_20261017.json} (2\,kB): the
certificate used in the proof of Proposition~\ref{prop:table}, with the
control on \cite[Table~3]{hb}. A few minutes.
\end{itemize}
Their SHA-256 values are:
{\footnotesize
\begin{align*}
&\text{P1}\quad\texttt{ab90cc7dee26ad872c0dff84e96a97bfea1d39838115fe6c277330f9a0b47690}\\
&\text{R1}\quad\texttt{e1e2f2b2b19973006484b5f79aa3373a440d981ba46003e28762879fd290341a}\\
&\text{P2}\quad\texttt{4f78eeab23aa25a676aa80eb61c24b4829d35789384497043c76c76f3139c3de}\\
&\text{R2}\quad\texttt{460a3478b9f0f63af370e3ffb3e884e584f7e7d123cc1cd9e4ac3bf9a1ba660b}\\
&\text{P3}\quad\texttt{94073d936706b5ca0ad6fa707a39e8566bee938c4cd4806a04af1122c5ef5935}\\
&\text{R3}\quad\texttt{69997f80743411bb952c727f1e8c25b710fe28b14d8a2f995967fb08383eff67}
\end{align*}}

Each record regenerates byte for byte with Python 3.13.16 on x86\_64 Linux,
where it was generated. With Python 3.10.12 on aarch64 Linux, R3 replays
byte for byte; R1 agrees in sections A and E--H and differs in 24 floats of
sections B--D (three identity differences below $10^{-15}$ in size and 21
measured deviations, by at most $4\cdot10^{-11}$ relative); R2 differs in
two floats of section A, at rounding level. A replay shows that the
arithmetic is reproducible, not that the floating evaluation is
outward-rounded.

\subsection{Same-model checks}

Two Claude subagents of the drafting session, with their own code, checked
the notes this manuscript integrates \cite{n01,n02}.
\begin{itemize}
\item For Sections~\ref{sec:saddle}--\ref{sec:numerics} (mpmath): the
identity of Lemma~\ref{lem:gauss} at four points to $10^{-30}$; the three
forms of Lemma~\ref{lem:shift} to $4\cdot10^{-31}$;
Lemmas~\ref{lem:euler}--\ref{lem:partial} in about 23\,000 randomized cases
with no violation; the identity in the proof of Proposition~\ref{prop:E}
to $4\cdot10^{-42}$ for random entire $\Theta$ and to $8\cdot10^{-28}$ for
the $\Theta$ of $\zeta$; Binet's formula to $10^{-29}$; and P1's $E$ at five
parameter points against its own implementation, with relative differences
at most $4\cdot10^{-5}$, two of them with P1's value below its own.
\item For Sections~\ref{sec:derivative} and~\ref{sec:strip} (mpmath at 34
digits): that nothing in \cite[Lemmas~2, 3 and 10, Proposition~11 and
Theorem~12]{hb} uses the half-plane except through $M(0)\le Q_0$; every
$k_i$ and gate in exact rationals; the polynomial condition on all 34 rows
with its own adaptive Bernstein scheme and a different compactification;
$C^*$ and $G^*$ on all rows, with the same smallest margins; the
comparison from $Q_0=1$, reaching $\frac1{25}$ at $t\approx0.16865$; the
identities of Section~\ref{sec:compare}; and the formulas of
Corollary~\ref{cor:derivative}, including the branch points of
Remark~\ref{rem:branch}.
\end{itemize}
Both checks found errors, corrected in the notes before this draft: status
words stronger than a grid evaluation supports; a misquoted contraction
constant; a missing invariant domain for the zeta saddle; the crude strip
bound allowed where $\mathcal Q_*$ is not justified (numerically harmless);
numbers that disagreed with the record; the bound
$\beta'\le\frac1{12}$ of Lemma~\ref{lem:stirling}, at first only sampled;
the domain of Corollary~\ref{cor:derivative}; an incomplete dependency
list; the word ``unconditional'' for the table's conclusion; and Decimal
powers in P3, now repeated directed multiplications.

Two further Claude subagents checked this manuscript against the notes,
the records and \cite{hb}, with their own code. They found no error in the
proofs of Sections~\ref{sec:setting}--\ref{sec:derivative}. Points they
raised, corrected here, include: conductor thresholds of
Table~\ref{tab:conductor} that the notes had rounded below the passing
values of the record; two
numbers of Section~\ref{sec:validation} that the notes had interchanged;
the claim that \cite[Proposition~7]{hb} shares the main term of
Theorem~\ref{thm:transport}; the hypotheses of
Corollary~\ref{cor:derivative} on the disk; and the dependency list of
Proposition~\ref{prop:table}. All these checks share the model family of
the drafting model and are not independent review.

\section*{Acknowledgement of LLM assistance}
This working manuscript was drafted for Edward Baker with substantial
assistance from Claude Opus 5.5 (Anthropic; configured identifier
\texttt{claude-opus-5-5}, Claude Cowork, reasoning effort max). The model
carried out the derivations, wrote the check programmes, ran the
literature search and drafted the text; same-model subagents made the
checks described in Appendix~\ref{app:repro}. These checks do not
constitute independent cross-model or specialist human review. The
preprint \cite{hb} on which Section~\ref{sec:strip} relies states that it
was prepared with ChatGPT assistance and has not been independently
refereed. This manuscript is unreviewed and does not imply the intended
author's endorsement of its mathematical claims.

\begin{thebibliography}{99}

\bibitem{debruijn}
N.~G. de Bruijn,
\emph{The roots of trigonometric integrals},
Duke Math. J. \textbf{17} (1950), 197--226.

\bibitem{dlmf}
NIST Digital Library of Mathematical Functions, \url{https://dlmf.nist.gov/},
Section~5.9(ii).

\bibitem{dobner}
A.~Dobner,
\emph{A proof of Newman's conjecture for the extended Selberg class},
Acta Arith. \textbf{201} (2021), 29--62; arXiv:2005.05142.

\bibitem{gomila}
J.~Gomila,
\emph{A computer-assisted proof of the bound $\Lambda\le0.1787854$ for the
de Bruijn--Newman constant}, unrefereed preprint, 20 August 2026,
repository \nolinkurl{judegomila/dbn-lambda-01787854-candidate-audit}
(commit a74738d). Read in part.

\bibitem{gordon}
S.~Gordon,
\emph{A computer-assisted upper bound for the de Bruijn--Newman constant},
unrefereed preprint, September 2026, repository
\nolinkurl{stefangordon/dbn-upper-bound} (commit 64c47f2). Read in part.

\bibitem{hb}
C.~Hakimoglu-Brown,
\emph{A conditional bound $\Lambda\le0.13963$ for the de Bruijn--Newman
constant}, arXiv:2610.12803v1, working draft of 7 October 2026,
unrefereed.

\bibitem{newman}
C.~M. Newman,
\emph{Fourier transforms with only real zeros},
Proc. Amer. Math. Soc. \textbf{61} (1976), 245--251.

\bibitem{p01}
Project~01 (Dirichlet--Newman families), Note~01,
\emph{Normalization and transfer contract}, 17 October 2026,
\nolinkurl{papers/investigations-3/01_dirichlet_newman_families/notes/01_NORMALIZATION_AND_TRANSFER_CONTRACT_CLAUDE_20261017.md},
in the repository of \cite{records}. Unreviewed.

\bibitem{n01}
Note~01, \emph{Height-$p$ transport and Lemma $A_p$}, 17 October 2026,
\nolinkurl{papers/investigations-3/03_effective_outer_remainders/notes/round-1/01_HEIGHT_P_TRANSPORT_AND_LEMMA_AP_CLAUDE_20261017.md}
(SHA-256 \texttt{f591b9e7\dots85a2}), in the repository of
\cite{records}.

\bibitem{n02}
Note~02, \emph{Hakimoglu-Brown's strip comparison against the height-$p$
floor, and a table from the unconditional strip}, 17 October 2026,
\nolinkurl{papers/investigations-3/03_effective_outer_remainders/notes/round-1/02_HAKIMOGLU_BROWN_COMPARISON_AND_UNCONDITIONAL_TABLE_CLAUDE_20261017.md}
(SHA-256 \texttt{3a95a8d3\dots77d5}), in the repository of
\cite{records}.

\bibitem{planat}
Planat, arXiv:2610.18214v1, 13 October 2026, unrefereed. Read in part
(HTML to Section~3).

\bibitem{platt2016}
D.~J. Platt,
\emph{Numerical computations concerning the GRH},
Math. Comp. \textbf{85} (2016), 3009--3027, doi:10.1090/mcom/3077;
arXiv:1305.3087.

\bibitem{platt2017}
D.~J. Platt,
\emph{Isolating some non-trivial zeros of zeta},
Math. Comp. \textbf{86} (2017), 2449--2467, doi:10.1090/mcom/3198.

\bibitem{platttrudgian}
D.~Platt and T.~Trudgian,
\emph{The Riemann hypothesis is true up to $3\cdot10^{12}$},
Bull. Lond. Math. Soc. \textbf{53} (2021), 792--797; arXiv:2004.09765.

\bibitem{polymath}
D.~H.~J. Polymath,
\emph{Effective approximation of heat flow evolution of the Riemann $\xi$
function, and a new upper bound for the de Bruijn--Newman constant},
Res. Math. Sci. \textbf{6} (2019), Paper No.~31; arXiv:1904.12438.

\bibitem{raymond}
Raymond,
\emph{Correlated translated probes and the bound $\Lambda<0.149396$},
Zenodo record 22867547, 21 September 2026, unrefereed. Read in part.

\bibitem{raymondcond}
Raymond, conditional bound $\Lambda<0.116398$ under the hypothesis that
$\zeta$ has no zero with $\re s>\frac78$, Zenodo record of 10 October 2026,
unrefereed. Read in part.

\bibitem{records}
Round-one code and records,
\nolinkurl{papers/investigations-3/03_effective_outer_remainders/numerics/round-1/},
repository \url{https://github.com/ebbaker/shifted-zeta-positivity}
(Appendix~\ref{app:repro}).

\bibitem{rodgerstao}
B.~Rodgers and T.~Tao,
\emph{The de Bruijn--Newman constant is non-negative},
Forum Math. Pi \textbf{8} (2020), e6.

\end{thebibliography}

\end{document}
